%% The Elicit installation manual — umbrella UC-001, "Consult the Installation Manual".
%%
%% Build it with docs/manual/build-manual.sh, which renders it through a TeX Live container
%% and stamps it with the platform version and build date (UC-002 BR-002). See README.md.
%%
%% Scope is provisioning a deployment (umbrella C-004), including the reports a site adds to
%% its surveys (UC-020, UC-021). Operating the console is the administrator's manual (Admin
%% UC-029). This document points at it and stops.

\documentclass[11pt, a4paper]{article}

\usepackage[a4paper, margin = 25mm, bottom = 28mm]{geometry}

\newcommand{\elicitProduct}{Elicit Platform}
\usepackage{elicit-brand}
\usepackage{hyperref}
\elicitAccent{elicitVerdigris}

\hypersetup{
  pdftitle = {Elicit — Installing the Platform},
  pdfauthor = {Elicit Software},
  pdfsubject = {Installing, configuring and customizing the Elicit platform},
  colorlinks = true,
  linkcolor = elicitInkDark,
  urlcolor = elicitInkDark,
}

\graphicspath{{./}}

%% The box in front of each line of the security checklist (UC-022 BR-004).
\newcommand{\elicitcheck}{\raisebox{-0.12em}{\setlength{\fboxsep}{0pt}\fbox{\rule{0pt}{0.72em}\hspace{0.72em}}}}

\begin{document}

%% ------------------------------------------------------------------ title page
\elicitTitlePage{Installing the Platform}{The Elicit installation manual}%
  {Survey \textperiodcentered\ Admin \textperiodcentered\ FHHS \textperiodcentered\ Pedigree.
   Author, the authoring tool, has an installation manual of its own.}%
  {This manual describes the platform version named above. It is a release artifact of the
   Elicit umbrella repository and is read before any Elicit service is running.\\[0.5em]
   Elicit Software \textperiodcentered\ PolyForm Noncommercial License 1.0.0}

%% ------------------------------------------------------------------ contents
\pagenumbering{roman}
\setcounter{page}{1}
\begin{elicitnote}[Who this manual is for]
Someone comfortable with PostgreSQL, a container runtime and an OIDC provider, who has not
seen Elicit before. Parts~I and II take a site from an empty machine to its first survey.
Part~III makes it the site's own: its brand, its languages and its reports, the last of which
is written for whoever builds a report service. Part~IV secures the site and checks the result, and the appendix is
the reference for every setting.

Running the console day to day is the administrator's manual. This one stops where that
begins.
\end{elicitnote}

\vspace{1.2em}
\elicitContents

\clearpage
\pagenumbering{arabic}
\setcounter{page}{1}

\elicitpart{Getting oriented}{What an Elicit site is made of, how its pieces depend on each
other, and what has to exist before you install any of them.}

%% ================================================================== 1
\section{What Elicit is made of}
\label{sec:topology}

Elicit is a modular survey platform: a handful of applications that share one database, one
identity provider and one brand directory. A \term{site} is one installation of them. This
chapter names the pieces and how they connect, so that you can decide which of them your site
runs before you install anything.

\begin{elicitdiagram}{A site. Respondents reach Survey and administrators reach Admin. The
optional report services are called by Survey and Admin, never by a browser.}
  \node[elicit person, minimum width = 2.6cm] (resp) at (0, 0) {Respondent};
  \node[elicit person, minimum width = 2.6cm] (adm) at (6.4, 0) {Administrator};
  \node[elicit app, minimum width = 2.6cm] (survey) at (0, -1.8) {Survey};
  \node[elicit app, minimum width = 2.6cm] (admin) at (6.4, -1.8) {Admin};
  \node[elicit service, minimum width = 2.3cm] (idp) at (10.0, -1.25) {Identity provider};
  \node[elicit service, minimum width = 2.3cm] (smtp) at (10.0, -2.35) {SMTP relay};
  \node[elicit database, minimum width = 4.4cm] (db) at (3.2, -3.6)
    {PostgreSQL\\[-1pt]{\scriptsize the \val{survey} schema and the reporting schemas}};
  \node[elicit app, minimum width = 2.6cm] (fhhs) at (0, -5.4)
    {FHHS\\[-1pt]{\scriptsize report service}};
  \node[elicit app, minimum width = 2.6cm] (ped) at (3.6, -5.4)
    {Pedigree\\[-1pt]{\scriptsize draws pedigrees}};
  \begin{scope}[on background layer]
    \node[elicit group, fit = (fhhs) (ped), inner ysep = 12pt] (opt) {};
  \end{scope}
  \node[elicit group label, anchor = south west] at (opt.north west) {optional};
  \draw[elicit flow] (resp) -- node[elicit label] {access code} (survey);
  \draw[elicit flow] (adm) -- node[elicit label] {signs in} (admin);
  \draw[elicit flow] (survey) -- (db);
  \draw[elicit flow] (admin) -- (db);
  \draw[elicit flow] (admin.east) -- (idp.west);
  \draw[elicit flow] (admin.east) -- (smtp.west);
  \draw[elicit flow] ([xshift = -0.9cm]survey.south) -- node[elicit label, pos = 0.38]
    {reports} ([xshift = -0.9cm]fhhs.north);
  \draw[elicit flow] (fhhs) -- (db);
  \draw[elicit flow] (fhhs) -- (ped);
\end{elicitdiagram}

\subsection{The modules}

\begin{plaintable}{tab:modules}{The modules, what each one is, and what it needs}{@{}>{\raggedright\arraybackslash}p{0.16\linewidth} >{\raggedright\arraybackslash}p{0.40\linewidth} >{\raggedright\arraybackslash}p{0.37\linewidth}@{}}{\thead{Module} & \thead{What it is} & \thead{What it needs}}
\ui{Survey} & The subject-facing application. A respondent enters an \term{access code} and
answers the survey. It owns the shared \val{survey} schema and creates it. &
The database. Required. \\
\addlinespace
\ui{Admin} & The administration console: departments, users, subjects, access codes,
message templates, reports, and applying survey definitions. &
The database, the identity provider, an SMTP relay. Required. \\
\addlinespace
\ui{FHHS} & The report service of the Family Health History Survey, and the reference
example of a custom report (Chapter~\ref{sec:fhhs}). Optional. &
The database, Pedigree. \\
\addlinespace
\ui{Pedigree} & Draws the family-tree diagrams FHHS asks for. Optional. & Nothing but the
service token. \\
\addlinespace
\ui{Author} & The authoring tool: designs surveys and exports the definition files Admin
applies. Optional, and installed from its own manual. &
A database of its own, the identity provider, a preview Survey. \\
\end{plaintable}

\begin{elicitnote}[Optional means optional]
A site that does not run the Family Health History Survey runs neither \ui{FHHS} nor
\ui{Pedigree}, and nothing outside Chapter~\ref{sec:fhhs} requires them. A site that only
delivers surveys, applying definition files written elsewhere, runs no \ui{Author}; the
authoring tool has an installation manual of its own.
\end{elicitnote}

\subsection{What they share}

\begin{itemize}
  \item \textbf{One PostgreSQL cluster.} A site's survey database carries the
        \val{survey} schema, the common reporting schema \val{surveyreport}, and one
        reporting schema per installed survey, which Survey creates itself. Survey, Admin
        and FHHS all connect to it.
  \item \textbf{One identity provider.} Admin authenticates against its own OIDC client.
        Survey authenticates no one: a respondent enters an access code, which is not a
        credential in the identity provider.
  \item \textbf{One brand directory,} mounted read-only into Survey and Admin (and into
        Author and its preview, where a site runs them).
  \item \textbf{One service token,} a secret of the site's own that every module presents when
        it calls another (Chapter~\ref{sec:service-token}).
\end{itemize}

Translations are optional. Survey and Admin carry English in the image. Latin American Spanish
(\val{es-419}) and Arabic (\val{ar}) come in a curated translations directory that
Elicit Software maintains beside each release, which a site mounts read-only, along with any
files of its own. A site that mounts nothing offers English only (Chapter~\ref{sec:i18n}).

\subsection{The order of the first start}

Survey creates and migrates the shared \val{survey} schema. Admin adds its own tables to that
schema, and FHHS adds its grants and indexes on top of both. So on a \emph{new} database:

\begin{enumerate}
  \item \ui{Survey} starts and becomes ready. Readiness is what tells you its migrations
        have finished.
  \item \ui{Admin} starts and becomes ready. Its migrations run against the schema Survey
        created.
  \item \ui{FHHS} starts. Started beside Admin rather than after it, the two modules' first
        migrations grant on the same tables at the same moment, and one of them fails.
  \item \ui{Pedigree} starts whenever you like: it has no database.
\end{enumerate}

\begin{elicitnote}[The order constrains the first start only]
It exists because of schema creation. Once the schema is migrated the modules may be started
in any order and may restart independently. The compose file enforces the order for you;
without it, the order is yours to keep (Chapter~\ref{sec:nocompose}).
\end{elicitnote}

\subsection{Ports}

Every application listens on port \val{8080} \emph{inside} its container. The host ports
below are what the shipped compose file publishes.

\begin{plaintable}{tab:ports}{Published host ports in the shipped compose file}{@{}>{\raggedright\arraybackslash}p{0.13\linewidth} >{\raggedright\arraybackslash}p{0.28\linewidth} >{\raggedright\arraybackslash}p{0.49\linewidth}@{}}{\thead{Host port} & \thead{Service} & \thead{Reached by}}
\val{8080} & Survey & Respondents, in a browser. \\
\val{8081} & Admin & Administrators, in a browser. \\
\val{8082} & FHHS & Survey and Admin, for reports. The shipped Family History definition
names FHHS at this port (Chapter~\ref{sec:fhhs}). \\
\val{8083} & Pedigree & You, for local testing. FHHS reaches it inside the compose network. \\
\val{8084} & Author & Survey authors, in a browser. See the Author installation manual. \\
\val{8085} & Author's preview Survey & Survey authors, in a browser. \\
\val{5452} & PostgreSQL (\val{5432} inside) & You, with a database tool. The applications reach
it inside the compose network. \\
\val{8180} & Keycloak (bundled) & Browsers. Admin and Author reach it inside the compose
network. \\
\val{8025} & Mailpit web interface (bundled) & You. \\
\val{1025} & Mailpit SMTP (bundled) & Admin. \\
\val{16686} & Jaeger user interface (bundled) & You. \\
\val{4317} & Jaeger OTLP/gRPC collector (bundled) & Every application. \\
\val{4318} & Jaeger OTLP over HTTP (bundled) & Nothing in the stack. \\
\end{plaintable}

\begin{elicitcaution}[Health checks address the container, not the host]
Each service's health check targets \val{http://localhost:8080/q/health/ready} \emph{inside}
the container. That is the same address for every application, whichever host port it is
published on. A health check pointed at the published host port is the most common
self-inflicted installation failure.
\end{elicitcaution}

\subsection{Metrics and traces}

Every application exposes Prometheus-format metrics at \opt{/q/metrics} and health at
\opt{/q/health/live} and \opt{/q/health/ready}, and exports OpenTelemetry traces to a
collector.

\begin{elicitnote}[There is no metrics stack in this repository]
Nothing in the shipped compose file scrapes \opt{/q/metrics}, and no Prometheus service runs.
The endpoint is there for your site's own scraper. The bundled Jaeger collects traces only, so
its \ui{Monitor} tab, which would read metrics, is inert.
\end{elicitnote}

The health and metrics endpoints answer without a sign-in, and logs and traces can name
respondents. Section~\ref{sec:security-observability} says who should be able to reach each on a
real site.

%% ================================================================== 2
\section{Before you start}
\label{sec:before}

Check each of these for the modules your site will run. A missing prerequisite is far cheaper
to find now than halfway through an installation.

\begin{description}
  \item[A PostgreSQL 17 cluster] in which you can create roles, databases and schemas.
  \item[A container runtime:] one with compose to evaluate Elicit
        (Chapter~\ref{sec:compose}), and for a real site the container platform it already
        runs (Chapter~\ref{sec:nocompose}).
        Elicit is released as container images only.
  \item[An OIDC provider] in which you can create a confidential client and three roles,
        reachable at \emph{one} address from both the browser and the application containers.
  \item[An SMTP relay] for Admin's invitations and reminders. Admin is the only module that
        sends mail.
  \item[A directory for the brand,] readable by every module that mounts it, and the same for
        translations if the site offers any language besides English (Chapter~\ref{sec:i18n}).
  \item[Free host ports] for the modules you will run (Table~\ref{tab:ports}).
\end{description}

\begin{elicitcaution}[One address for the identity provider]
The issuer address you configure must resolve to the same provider \emph{from the browser}
and \emph{from inside the container}. When it does not, sign-in fails in a way that looks
like an application fault: the browser is redirected somewhere it cannot reach, or, on a host
that maps the provider's name to itself, lands on whichever service listens on that port. The
bundled compose file does not rely on one name for both: Keycloak is given its browser address
(\opt{KC_HOSTNAME}\val{=http://localhost:8180}), which it puts in every sign-in redirect and
names as issuer, and the applications call it at its service name,
\val{http://keycloak:8080}. No hosts entry is needed on the evaluation host.
\end{elicitcaution}

The repository bundles a database, an identity provider, a mail catcher, an SFTP server and a
trace collector so that an evaluation works in one command (Chapter~\ref{sec:compose}). They
are for development and evaluation: a real site uses none of them, and none of the documented
passwords they ship with (Chapter~\ref{sec:secrets}).

\elicitpart{Installing}{From an empty database cluster to a site with its first survey: the
database, the identity provider, the secrets, the service token, the bundled stack for
evaluating Elicit and the
installation of a real site, the
first sign-in, and applying a survey definition.}

%% ================================================================== 3
\section{The database}
\label{sec:database}

Elicit creates no database roles: every role must exist before any module starts. Nor do you
create any tables. You create two empty schemas, and everything inside them comes from the
modules' migrations on their first start.

\subsection{The roles}

\begin{plaintable}{tab:roles}{The database roles Elicit expects}{@{}>{\raggedright\arraybackslash}p{0.24\linewidth} >{\raggedright\arraybackslash}p{0.14\linewidth} >{\raggedright\arraybackslash}p{0.56\linewidth}@{}}{\thead{Role} & \thead{Kind} & \thead{Used by}}
\val{dbowner} & group & Holds \val{CREATE} on the Elicit database. The schema owner is a member
and has it through the group. \\
\val{dbuser} & group & Holds the connect and usage privileges every application role
inherits. \\
\addlinespace
\val{elicit_owner} & login & The \emph{schema owner}. Every module's Flyway migrations run as
this role, and Survey builds each survey's reporting schema as it. \\
\val{survey_user} & login & The application connection of Survey and of FHHS. \\
\val{surveyadmin_user} & login & Admin's application connection. \\
\val{surveyreport_user} & login & Reads the reporting schemas, for a report service of the
site's own (Chapter~\ref{sec:reports}). \\
\end{plaintable}

Run this once for the cluster, as a superuser of the cluster, substituting a password of your
own for each placeholder. None of the roles it creates is a superuser.

\begin{elicitsql}
CREATE ROLE dbowner WITH NOLOGIN NOSUPERUSER INHERIT
  NOCREATEDB NOCREATEROLE NOREPLICATION;

CREATE ROLE dbuser WITH NOLOGIN NOSUPERUSER INHERIT
  NOCREATEDB NOCREATEROLE NOREPLICATION;

CREATE ROLE elicit_owner WITH LOGIN NOSUPERUSER INHERIT
  NOCREATEDB NOCREATEROLE NOREPLICATION
  ENCRYPTED PASSWORD '<owner password>';
GRANT dbowner TO elicit_owner;

CREATE ROLE survey_user WITH LOGIN NOSUPERUSER INHERIT
  NOCREATEDB NOCREATEROLE NOREPLICATION
  ENCRYPTED PASSWORD '<survey password>';
GRANT dbuser TO survey_user;

CREATE ROLE surveyadmin_user WITH LOGIN NOSUPERUSER INHERIT
  NOCREATEDB NOCREATEROLE NOREPLICATION
  ENCRYPTED PASSWORD '<admin password>';
GRANT dbuser TO surveyadmin_user;

CREATE ROLE surveyreport_user WITH LOGIN NOSUPERUSER INHERIT
  NOCREATEDB NOCREATEROLE NOREPLICATION
  ENCRYPTED PASSWORD '<report password>';
GRANT dbuser TO surveyreport_user;
\end{elicitsql}

\begin{elicitnote}[What the schema owner can do]
\val{elicit_owner} is not a superuser, and the modules' migrations are tested with it as the
script creates it. It can:
\begin{itemize}
  \item create and alter objects in the \val{survey} and \val{surveyreport} schemas, which it
        owns;
  \item create a schema for each installed survey, through its membership of \val{dbowner}
        and the \val{CREATE ON DATABASE} grant below, which is required rather than a
        courtesy;
  \item own, alter, rename and drop what it creates there.
\end{itemize}
It cannot create roles or databases, read another database of the cluster, or reach the
server's files. A cluster whose roles were created by an earlier release's script, or by
version 1.0.0 of the bundled database image, still has \val{elicit_owner} and \val{dbowner} as
superusers: \file{DeploymentScript.md} gives the three statements that narrow them.
\end{elicitnote}

\subsection{The site's survey database}

Create a database named \val{survey}, then run this inside it. \val{surveyreport} is the
common reporting schema: it holds only the date and status dimensions every survey shares.
Each survey's own star schema, named \val{report_...} after the survey and recorded on
\val{survey.surveys.report_schema}, is created by Survey when the survey is first built, not
here; the \val{CREATE ON DATABASE} grant is what allows it.

\begin{elicitsql}
GRANT CONNECT, CREATE ON DATABASE survey TO dbowner;
GRANT CONNECT ON DATABASE survey TO dbuser;

GRANT USAGE, CREATE ON SCHEMA public TO dbowner;
GRANT USAGE ON SCHEMA public TO dbuser;

CREATE SCHEMA survey AUTHORIZATION elicit_owner;
GRANT USAGE ON SCHEMA survey TO survey_user;
GRANT USAGE ON SCHEMA survey TO surveyadmin_user;
GRANT USAGE ON SCHEMA survey TO surveyreport_user;

CREATE SCHEMA surveyreport AUTHORIZATION elicit_owner;
GRANT USAGE ON SCHEMA surveyreport TO survey_user;
GRANT USAGE ON SCHEMA surveyreport TO surveyadmin_user;
GRANT USAGE ON SCHEMA surveyreport TO surveyreport_user;

GRANT USAGE, CREATE ON SCHEMA surveyreport TO dbowner;
ALTER DEFAULT PRIVILEGES IN SCHEMA surveyreport
  GRANT SELECT, INSERT, UPDATE, DELETE ON TABLES TO dbowner;
ALTER DEFAULT PRIVILEGES IN SCHEMA surveyreport
  GRANT USAGE ON SEQUENCES TO dbowner;
\end{elicitsql}

Both schemas are now empty. That is correct: every table, sequence, index and grant inside
them is created by a module's migrations the first time it starts.

\begin{elicitnote}[The bundled database image does all of this for you]
\file{elicitsoftware/elicit_db} runs the role script and the schema script when its data
directory is first initialized. The compose file names version 3.0.0-rc.1 of the image, the first
whose roles are the ones above. An initialization script runs \emph{only} against an empty
data directory, so a cluster that already exists is yours to prepare with the two scripts
above. The image's scripts carry the documented development passwords listed in
Chapter~\ref{sec:secrets}.
\end{elicitnote}

%% ================================================================== 4
\section{The identity provider}
\label{sec:oidc}

Admin authenticates against a confidential OIDC client. Survey has no client: a respondent
enters an access code, which is issued by Admin and is not a credential in your identity
provider.

Elicit does not depend on Keycloak. The realm the repository ships (Section~\ref{sec:oidc-example})
is an example, and any OpenID Connect provider will do: your organization's own, a cloud
directory, or a Keycloak you run yourself. What the console needs from a sign-in is a
\emph{principal with a user name} (Section~\ref{sec:oidc-username}); everything else it
decides for itself. It never reads a Keycloak-specific claim and never calls Keycloak's
administration API.

The same holds beyond OIDC. The console reads only the signed-in principal's name and its
roles, so any authentication method that yields a principal with a user name will work, but
the principal has to be established \emph{inside} Admin. The published image establishes it
with OIDC only, and it ignores a user name a reverse proxy passes along: a web server that signs
people in by Basic, LDAP or SAML and forwards \val{REMOTE_USER} or a header gets them to the
console's front door, where Admin sends them to the identity provider all the same. Choosing
\val{DATABASE} for the roles (Section~\ref{sec:oidc-authorization}) does not change this, since
it decides roles only after someone is signed in. A proxy that is itself an OIDC client can
forward the provider's access token as \val{Authorization: Bearer}, which Admin accepts; any
other method means building Admin with an authentication mechanism for it.

\subsection{The client}

Create one client in the provider's realm for Admin. Survey needs none.

\begin{plaintable}{tab:clients}{The OIDC client}{@{}>{\raggedright\arraybackslash}p{0.22\linewidth} >{\raggedright\arraybackslash}p{0.20\linewidth} >{\raggedright\arraybackslash}p{0.52\linewidth}@{}}{\thead{Client} & \thead{Type} & \thead{Redirect}}
\val{elicit-admin} & confidential & The console's address. On a real site that is its
public \val{https} address and nothing wider, for example \val{https://console.example.org/*};
the evaluation stack uses \val{http://localhost:8081/*}. \\
\end{plaintable}

The client needs a secret, which reaches the application as \opt{OIDC_CLIENT_SECRET}. Admin
accepts browser sign-ins and bearer tokens on the same client.

\screenshot{01-keycloak-clients}{The bundled realm's clients. A site creates the console's
client in its own provider, under whatever name it configures Admin with.}

\screenshot{02-keycloak-client}{The console's client. The redirect URI must cover the address
a browser reaches the console at, which is not necessarily the address the container knows.}

\subsection{The user name}
\label{sec:oidc-username}

The console identifies a person by the name of their principal and nothing else. After every
sign-in it looks that name up among its own user records (\val{survey.users.username}), and
only an active record with that exact name can use the console: it carries the person's
departments, and without one the console logs \val{No active user found for authenticated
principal} and offers nothing. This holds in both authorization modes
(Section~\ref{sec:oidc-authorization}). A machine client's name is matched the same way,
against its import client record.

With OIDC the name is taken from a claim of the token: \val{upn} if the token has one,
otherwise \val{preferred_username}, otherwise \val{sub}. Keycloak puts the account's user name
in \val{preferred_username}. A provider whose user name is in another claim, such as
\val{email}, is followed by setting \val{quarkus.oidc.token.principal-claim} on Admin
(environment: \val{QUARKUS_OIDC_TOKEN_PRINCIPAL_CLAIM}). Choose a claim whose value does not
change while the person keeps their account, and create each user record in the console with
that value as its username.

\subsection{The roles}

\begin{plaintable}{tab:roles-oidc}{The roles the applications check}{@{}>{\raggedright\arraybackslash}p{0.20\linewidth} >{\raggedright\arraybackslash}p{0.18\linewidth} >{\raggedright\arraybackslash}p{0.56\linewidth}@{}}{\thead{Role} & \thead{On client} & \thead{What it permits}}
\val{elicit_admin} & \val{elicit-admin} & The console's administrative screens. \\
\addlinespace
\val{elicit_user} & \val{elicit-admin} & The console's day-to-day screens, without the
administrative ones. \\
\addlinespace
\val{elicit_importer} & \val{elicit-admin} & The subject-import endpoints of the
integration API, for a machine client presenting a bearer token, and nothing else: no
person holds it and neither of the other two roles implies it. The role lets the client
call; the departments it registers into come from its import client record in the console. \\
\end{plaintable}

At least one account must hold \val{elicit_admin} before the first sign-in
(Chapter~\ref{sec:firstsignin}), or nobody can create the first department.

\screenshot{03-keycloak-roles}{The roles on the console's client. A site creates the same
ones in its own provider.}

\subsection{The bundled example}
\label{sec:oidc-example}

The repository ships a Keycloak realm at \file{keycloak/elicit-realm.json}, carrying the
console's client and its roles, the integration client \val{import-service}, Author's client,
and the accounts \val{admin}, \val{user}, \val{author} and \val{etl}. It is a worked example
and an evaluation provider, not a production one. Everything in the file is published with the
repository: the secret of the console's client and of Author's is the literal string
\val{secret}, the integration client's is \val{importer}, every account's password is
documented, and the realm's signing keys are the same in every copy. Keycloak's own bootstrap
administrator, in its \val{master} realm, is \val{admin}/\val{admin}.
The integration client has no import client record in a new database, so the registration API
answers it 403 until one is created under \ui{Admin} $>$ \ui{Import Clients}, named
\val{service-account-import-service}, with a department. The account \val{etl} holds
\val{elicit_importer} as a person would and is refused whatever is entered: the role is an
import client's alone.

A real site does not import this file or change its secrets one by one. It creates the client
and the roles above in its own identity provider, which issues secrets and keys of its own.

\subsection{Where the roles come from}
\label{sec:oidc-authorization}

Authentication, proving who someone is, is always the identity provider's. Authorization,
deciding which of the roles in Table~\ref{tab:roles-oidc} they hold, can come from either the
provider or the console's own database, chosen by \opt{elicit.authorization.mode} on Admin.

\begin{description}
  \item[\val{OIDC}, the default.] The roles arrive in the token, and the provider is the only
        place they are granted or withdrawn. \ui{Users} in the console then assigns
        departments only. The roles are read from the token's \val{groups} claim or, from
        Keycloak, from the realm's and the console client's role lists; for a browser sign-in
        the ID token is read. A provider that puts them elsewhere is followed with
        \val{quarkus.oidc.roles.role-claim-path} (the claim holding them) and
        \val{quarkus.oidc.roles.source} (\val{idtoken}, \val{accesstoken} or \val{userinfo}).
        Many providers can also be told to emit the role names as \val{groups}.
  \item[\val{DATABASE}.] The provider only authenticates and need not know Elicit's roles at
        all. An administrator assigns them in the console's user editor, which shows a role
        section in this mode, and they are kept in \val{survey.user_roles}. The table is
        consulted only for a principal whose token carries none of Elicit's roles: a role the
        provider does supply still wins, and a user record that is inactive grants nothing.
\end{description}

Neither mode turns OIDC off or changes sign-in, and in both the principal's name must match an
active user record (Section~\ref{sec:oidc-username}). Choose \val{DATABASE} when the people who
run the console cannot have roles created in the organization's provider; otherwise keep the
roles at the provider (Section~\ref{sec:security-identity}).

%% ================================================================== 5
\section{Replacing the development secrets}
\label{sec:secrets}

The repository and its images ship with working credentials so that the evaluation stack
starts in one command (Chapter~\ref{sec:compose}). None of them may be in use on a real site:
a site that was prepared from the scripts and examples in this manual checks every row of the
table below.

\begin{elicitsecurity}[Every one of these is published]
Each value below is in this repository and in this manual. A site that leaves one as shipped
is running with a credential anyone can read.
\end{elicitsecurity}

\begin{plaintable}{tab:secrets}{Development credentials that must be replaced}{@{}>{\raggedright\arraybackslash}p{0.29\linewidth} >{\raggedright\arraybackslash}p{0.21\linewidth} >{\raggedright\arraybackslash}p{0.43\linewidth}@{}}{\thead{What} & \thead{Shipped value} & \thead{Where it comes from}}
Schema owner password & \val{SURVEYPW} & The database image's role script;
\opt{ELICIT_OWNER_DB_PASSWORD} in the compose file, for every module. \\
\addlinespace
Survey's and FHHS's connection password & \val{SURVEYPW} & The database image's role script;
\opt{ELICIT_DB_PASSWORD} on the \val{survey} and \val{fhhs} services in the compose file. \\
\addlinespace
Admin's connection password & \val{SAOWNERPW} & The database image's role script;
\opt{ELICIT_DB_PASSWORD} in the compose file. \\
\addlinespace
Reporting password & \val{SRUSERPW} & The database image's role script. No shipped module
presents it; a report service of the site's own does (Chapter~\ref{sec:reports}). \\
\addlinespace
OIDC client secret, the console & \val{secret} & \opt{OIDC_CLIENT_SECRET} in the compose
file, and the \val{elicit-admin} client of the bundled Keycloak realm. \\
\addlinespace
OIDC client secret, integration client & \val{importer} & The \val{import-service} client of
the bundled Keycloak realm, used by a machine client presenting a bearer token. \\
\addlinespace
Database superuser & \val{survey}/\val{admin} & The bundled database service. \\
\addlinespace
Keycloak bootstrap administrator & \val{admin}/\val{admin} & The bundled identity provider's
\val{master} realm. \\
\addlinespace
Everything else in the bundled realm & as in the file & Author's client secret, the accounts
\val{admin}, \val{user}, \val{author} and \val{etl}, and the realm's signing keys. A real site
replaces none of them one by one: it uses its own provider and never imports the realm
(Chapter~\ref{sec:oidc}). \\
\addlinespace
Author's connection password & \val{SURVEYPW} & \opt{ELICIT_DB_PASSWORD} on the \val{author}
and \val{author-survey} services in the compose file. Author connects as \val{survey_user}, so
this is Survey's password again; the Author installation manual covers it. \\
\addlinespace
Service token & \val{elicit-dev-service-token-not-for-a-real-site} &
\opt{ELICIT_SERVICE_TOKEN} on every application service in the compose file, unless the
variable is set in the environment or in \file{.env}. A module running with it warns about it
in its log at every start. A site generates a value of its own
(Chapter~\ref{sec:service-token}). \\
\addlinespace
SFTP account & \val{fhhs_user}/\val{pass} & The bundled SFTP service, for FHHS report
delivery. \\
\addlinespace
SFTP key pair & the files in \file{sftp_ssh/} & A fixture for the bundled SFTP service, which
accepts its public half. A site generates a key pair of its own
(Section~\ref{sec:fhhs-sftp}). \\
\end{plaintable}

\begin{elicitcaution}[Every module needs both variables]
Survey, Admin and FHHS read their database passwords from \opt{ELICIT_DB_PASSWORD} and
\opt{ELICIT_OWNER_DB_PASSWORD}, and none of them ships a value, so no module ever connects with
a published password by default. A module started without them does not stop at once, though.
The database refuses it, and the module logs \val{no password was provided}, reports itself
not ready, and retries \opt{quarkus.flyway.owner.connect-retries} times with waits that double
up to two minutes. Only then does it exit with \val{Failed to start application}: about eight
minutes after it started, at the shipped ten retries. A module that sits unhealthy with that
message in its log is missing these variables.

Before 3.0.0, Survey and FHHS carried the literal \val{SURVEYPW}, so a site upgrading from an
earlier release must now set both variables on them, or \opt{quarkus.datasource.password} and
\opt{quarkus.datasource.owner.password}, which take precedence.
\end{elicitcaution}

\begin{elicitnote}[The two Admin variables read backwards, and are correct]
On the \val{admin} service the compose file sets \opt{ELICIT_DB_PASSWORD}\val{=SAOWNERPW} and
\opt{ELICIT_OWNER_DB_PASSWORD}\val{=SURVEYPW}. That is right: Admin's application connection
is \val{surveyadmin_user}, whose password is \val{SAOWNERPW}, and its migrating connection is
\val{elicit_owner}, whose password is \val{SURVEYPW}. The names simply do not read that way.
Do not swap them.
\end{elicitnote}

Two more things a new installation ships with are not secrets but are just as public: the
seeded accounts \val{admin} and \val{user}, which Admin warns about on every screen until
they are renamed (Chapter~\ref{sec:firstsignin}).

\subsection{Changing a password on a cluster that exists}
\label{sec:secrets-change}

The database image's role script runs only when its data directory is first initialized, so
editing it changes nothing on a cluster that already exists. There, and on a cluster of your
own, change a role's password with a statement, as a role that may alter it. Give each role a
different password, and write none of them into a file you keep under version control.

\begin{elicitsql}
ALTER ROLE elicit_owner      PASSWORD '<new owner password>';
ALTER ROLE survey_user       PASSWORD '<new survey password>';
ALTER ROLE surveyadmin_user  PASSWORD '<new admin password>';
ALTER ROLE surveyreport_user PASSWORD '<new report password>';
\end{elicitsql}

Then give the new value to every module that presents it, and restart those modules: Survey
first, then Admin, then FHHS.

\begin{plaintable}{tab:secrets-present}{Which module presents each role's password}{@{}>{\raggedright\arraybackslash}p{0.26\linewidth} >{\raggedright\arraybackslash}p{0.68\linewidth}@{}}{\thead{Role} & \thead{Presented as}}
\val{elicit_owner} & \opt{ELICIT_OWNER_DB_PASSWORD} on Survey, Admin and FHHS, and on Author's
preview Survey where a site runs it. \\
\addlinespace
\val{survey_user} & \opt{ELICIT_DB_PASSWORD} on Survey and FHHS, and on Author and its preview
Survey where a site runs them. \\
\addlinespace
\val{surveyadmin_user} & \opt{ELICIT_DB_PASSWORD} on Admin. \\
\addlinespace
\val{surveyreport_user} & Whatever a report service of the site's own reads it from
(Chapter~\ref{sec:reports}). \\
\end{plaintable}

A module still holding the old password is refused by the database and reports itself not
ready, as the caution above describes; nothing is lost, and giving it the new value is the whole
fix. The same two-place rule holds for every other secret (Section~\ref{sec:security-secrets}):
an OIDC client secret is changed in the identity provider and in \opt{OIDC_CLIENT_SECRET}, an
SFTP key on the server and in FHHS's settings (Table~\ref{tab:fhhs}), and a post-survey action's
bearer token at the receiving service and on Survey (Section~\ref{sec:psa}). The service token
has no issuer: it is changed on every module at once (Section~\ref{sec:service-token-rotate}).

%% ================================================================== 6
\section{The service token}
\label{sec:service-token}

The modules of a site call each other. Admin asks Survey to build a reporting schema, Survey and
Admin ask FHHS for a report, FHHS asks Pedigree for a drawing. Every such call carries one secret
value, the \term{service token}, and the module that is called checks it. Every module of a
site holds the same value, the value belongs to the site, and no image is shipped with one.

A caller presents the token in a request header, as
\val{Authorization: Bearer <token>}. It has nothing to do with a respondent's access
code, and it never reaches a browser. A module checks it on the paths it serves for other
modules and on no others:

\begin{plaintable}{tab:service-token-paths}{Where each module requires the service token}{@{}>{\raggedright\arraybackslash}p{0.16\linewidth} >{\raggedright\arraybackslash}p{0.41\linewidth} >{\raggedright\arraybackslash}p{0.35\linewidth}@{}}{\thead{Module} & \thead{Requires it on} & \thead{Answers without it}}
Survey & Every path under \opt{/api} except \opt{/api/pdf/download}: building, renaming and
dropping reporting schemas, and the respondent report. & Every page a respondent uses;
\opt{/api/pdf/download}, which the respondent's own session decides; and \opt{/q}. \\
\addlinespace
FHHS & Every path outside \opt{/q}: its three reports and its post-survey action. &
\opt{/q}. \\
\addlinespace
Pedigree & Every path but \opt{/healthz}. & \opt{/healthz}. \\
\addlinespace
Admin & Nothing. No module calls Admin; it holds the token to present it. & Its pages ask
for a sign-in, and \opt{/q} does not. \\
\end{plaintable}

\noindent On those paths a module gives one of three answers:

\begin{plaintable}{tab:service-token-answers}{What a protected path answers}{@{}>{\raggedright\arraybackslash}p{0.36\linewidth} >{\raggedright\arraybackslash}p{0.58\linewidth}@{}}{\thead{The request} & \thead{The answer}}
Carries the site's token. & It is served. \\
\addlinespace
Carries no token, or a wrong one. & \val{401}. The answer does not say which of the two it
was. \\
\addlinespace
Reaches a module that has no token itself. & \val{503}, with the message
\val{service token not configured}. A module without a token never serves these paths
unchecked. \\
\end{plaintable}

The health endpoints stay open, so a container platform probes a module without holding
the token (Section~\ref{sec:security-observability}).

\subsection{Generating it}

\begin{elicitcmd}
openssl rand -hex 32
\end{elicitcmd}

That prints 64 hexadecimal characters, which need no quoting in a shell, a properties file or a
YAML file. Any value of at least 32 characters is accepted, and a shorter one is refused.
Generate one value per site. A second site has its own, and so does an Author installation run
apart from a site.

Keep it where the site keeps its database passwords and supply it the same way: as an
environment variable from your platform's secret store, never in a file under version control
(Section~\ref{sec:security-secrets}).

\begin{elicitnote}[The evaluation stack has one already]
The shipped compose file gives every module the value
\val{elicit-dev-service-token-not-for-a-real-site}, so that the stack starts in one command
(Chapter~\ref{sec:compose}). The value is published (Table~\ref{tab:secrets}), and a module
running with it logs a warning at every start. Setting \opt{ELICIT_SERVICE_TOKEN} in the
environment, or in the \file{.env} file beside the compose file, replaces it for every service
at once.
\end{elicitnote}

\subsection{Which modules read it}
\label{sec:service-token-modules}

Every module the site runs gets the same value.

\begin{plaintable}{tab:service-token-modules}{How each module is given the service token}{@{}>{\raggedright\arraybackslash}p{0.26\linewidth} >{\raggedright\arraybackslash}p{0.68\linewidth}@{}}{\thead{Module} & \thead{Reads it from}}
Survey, Admin, FHHS & The environment variable \opt{ELICIT_SERVICE_TOKEN}, which supplies the
setting \opt{elicit.service.token} (Tables~\ref{tab:survey}, \ref{tab:admin}
and~\ref{tab:fhhs}). \\
\addlinespace
Pedigree & The environment variable \opt{ELICIT_SERVICE_TOKEN}. It is the only setting Pedigree
has. \\
\addlinespace
Author and its preview Survey & The same variable. The Author installation manual covers
them. \\
\end{plaintable}

\noindent What a module does when the value is missing or unusable:

\begin{description}
  \item[Not set.] The module starts, logs one warning that names the variable, and reports
        itself \emph{not ready}: \opt{/q/health/ready} shows a check named \val{service-token}
        as down, with the reason that \opt{ELICIT_SERVICE_TOKEN} is not set. Its protected
        paths answer \val{503}, and it makes no call that needs the token. The pages it serves
        to people are not affected. Pedigree has no readiness check to fail: \opt{/healthz}
        still answers, and \opt{/svg} answers \val{503}.
  \item[Shorter than 32 characters.] The module refuses to start. Its log names the variable
        and the length required, and never the value.
  \item[The development value.] The module works, and warns about the value in its log at
        every start. Admin's \ui{System} overview lists it among the setup work still
        outstanding.
\end{description}

No module writes the token to a log, a trace or an error message. Where a module reports on
its configuration, it says whether the token is set and nothing more.

\subsection{Where the token is sent}
\label{sec:service-token-origins}

A module sends the token to two kinds of address and to no other.

\begin{description}
  \item[A module it is configured with.] Admin sends it to Survey at \opt{elicit.survey.url},
        and FHHS to Pedigree at \opt{pedigree.url}. The site named those addresses in the
        module's own configuration, so there is nothing more to list.
  \item[A trusted origin.] The addresses of a survey's report services and post-survey
        actions arrive in its definition file (Chapter~\ref{sec:reports}). A module does not
        send the site's token to an address because a file named it: a mistaken definition, or
        a hostile one, would collect it. It sends the token only when the address's
        \term{origin} is one the site has listed in \opt{elicit.service.trusted-origins}.
\end{description}

An origin is a scheme, a host and a port, such as \val{http://fhhs.internal:8080}; the port
is \val{80} or \val{443} when none is written. The setting is a comma-separated list of them,
given as the environment variable \opt{ELICIT_SERVICE_TRUSTED_ORIGINS}. Origins are compared
exactly. There are no wildcards and no partial matches, and another port, or another name for
the same host, is another origin. An entry that carries a path, a query or a user name stops
the module at startup, with a message that names the entry.

\subsubsection*{The Family History reports, as the example}

The Family History definition lists FHHS's three reports and its post-survey action at one
origin (Chapter~\ref{sec:fhhs}). Three modules call those addresses: Survey, for the report
page and the action; Admin, for \ui{Print reports}; and FHHS itself, which asks its own report
addresses for the reports when it builds the document it files. So a site whose definition
names FHHS at \val{http://fhhs.internal:8080} sets, on Survey, Admin and FHHS:

\begin{elicitcmd}
ELICIT_SERVICE_TRUSTED_ORIGINS=http://fhhs.internal:8080
\end{elicitcmd}

The shipped compose file sets \val{http://host.docker.internal:8082} on those three services,
which is the origin the shipped definition names. A site that changes the report addresses
lists the origin it changed them to.

An address whose origin is not listed is still called, without the token. A module that
requires the token, as FHHS does, answers that call \val{401}, and the report shows an error
card in its place (Chapter~\ref{sec:trouble}).

\begin{elicitsecurity}[Listing an origin hands it the site's token]
Every call to a listed origin carries the token, and whoever holds the token can call every
protected path of every module. List only services the site runs itself and trusts as it
trusts Survey. A report service of the site's own does not have to be listed: left off, it is
called without a credential, and the network is its lock
(Section~\ref{sec:reports-network}).
\end{elicitsecurity}

Three more rules follow from the same caution:

\begin{itemize}
  \item A trusted origin on the private network may be \val{http}. One on another host
        should be \val{https}, because the token travels in the request as it is.
  \item A post-survey action that has a bearer token of its own (Section~\ref{sec:psa}) is
        sent that token and never the service token, whatever its address. One without is
        sent the service token if its origin is trusted, and no credential otherwise.
  \item A call that carries the token follows no redirect, and the reachability checks on
        Admin's \ui{Connections} page and in Survey's startup diagnostics never send it. To
        those checks, a protected address that answers \val{401} is reachable.
\end{itemize}

\subsection{Rotating it}
\label{sec:service-token-rotate}

Change the token when someone who knew it leaves, after any incident, and whenever your policy
says.

\begin{enumerate}
  \item Generate a new value.
  \item Replace the old value for every module of the site.
  \item Restart all of them together, not one at a time.
  \item Confirm with the checks in Chapter~\ref{sec:verify}.
\end{enumerate}

There is one token, and no period in which an old one is accepted beside a new one: the old
value stops working when the last module restarts. Between the first restart and the last, a
call from a module on one value to a module on the other is answered \val{401}. A definition
applied in that time reports that the reporting schema was not rebuilt, a report shows its
error card, and a post-survey action is retried on Survey's schedule (Section~\ref{sec:psa}).
Nothing is lost, and applying the definition again rebuilds the schema.

\subsection{What it does not do}
\label{sec:service-token-limits}

The token is one shared secret. It closes the paths between modules to anyone who is not a
module of the site, and it is worth knowing what it leaves as it was.

\begin{itemize}
  \item \textbf{It says that a module of the site is calling, not on whose behalf.} A caller
        that holds it may ask for any respondent's report, and for a build or a rename of any
        survey's reporting schema. Which respondent a person may see is decided by the
        application before it calls: Survey asks only for the respondent of the current
        session, and Admin asks after its own role and department checks.
  \item \textbf{Every holder can call every module.} Pedigree needs the token only to check
        the calls it receives, yet the value it holds is one FHHS and Survey would accept. The
        same is true of every service on the trusted list.
  \item \textbf{The module that faces the Internet holds it.} Survey both checks the token and
        presents it, so whoever could read Survey's environment would have it, as they would
        have Survey's database password.
  \item \textbf{It does not identify the caller.} A module that accepts a call knows that the
        caller held the token, not which module it was.
  \item \textbf{It is not encrypted by itself.} Over plain \val{http} it travels as written,
        as it does inside the compose network.
  \item \textbf{It changes for every module at once.}
\end{itemize}

So the token is a second lock and not a replacement for the first. Report services stay on
the private network (Section~\ref{sec:reports-network}), and the reverse proxy still refuses
the paths that are not for the people it serves (Section~\ref{sec:security-exposure}).

%% ================================================================== 7
\section{Evaluating with Docker Compose}
\label{sec:compose}

\file{docker-compose.yml}, in the umbrella repository and attached to each release, is a
complete Elicit site in one file that starts in one command: every application, and a database,
an identity provider, a mail catcher, an SFTP server and a trace collector for them to use. It
needs nothing beside it: no checkout of the repository, no brand or translations directory and
no data directory. It is how Elicit is developed, how its end-to-end tests run, and the quickest
way to see what the platform does.

\begin{elicitsecurity}[For development and evaluation, not for a real site]
The stack is built to be easy to start, not to be locked down
(Section~\ref{sec:compose-open}), and it is not the starting point of a production
installation. Enter no real respondent's data into it. A site that will hold real answers is
installed against its own database, identity provider and network
(Chapter~\ref{sec:nocompose}) and secured as Chapter~\ref{sec:security} describes;
Section~\ref{sec:compose-onward} says what an evaluation carries over to it.
\end{elicitsecurity}

\subsection{What the file defines}

Six application services and five supporting services:

\begin{itemize}
  \item the site: \val{survey}, \val{admin}, and the Family History reports \val{fhhs} and
        \val{pedigree} (Chapter~\ref{sec:fhhs});
  \item the authoring tool: \val{author} and its preview, \val{author-survey}, which the
        Author installation manual covers;
  \item what they depend on: \val{db}, \val{keycloak}, \val{mailpit}, \val{sftpServer} and
        \val{jaeger}.
\end{itemize}

Every Elicit image is \file{elicitsoftware/<name>} at one tag,
\val{${ELICIT_TAG:-<release>}}, so the applications, the database (\file{elicitsoftware/elicit_db}) and the bundled identity
provider (\file{elicitsoftware/elicit_keycloak}, Keycloak with the example realm built in)
always run as a matched set from one release. The tag defaults to the release the file came with, never to
\val{latest}; you run another release by naming its version (Section~\ref{sec:compose-up}).

The file mounts nothing from the host. What has to outlive a container is kept in named
volumes, which \val{docker compose down} keeps and \val{docker compose down -v} deletes:

\begin{plaintable}{tab:volumes}{Where the shipped compose file keeps its data}{@{\hspace{4pt}}>{\raggedright\arraybackslash}p{0.22\linewidth} >{\raggedright\arraybackslash}p{0.2\linewidth} >{\raggedright\arraybackslash}p{0.48\linewidth}@{\hspace{4pt}}}{\thead{Volume} & \thead{Service} & \thead{What it holds}}
\val{pgdata} & \val{db} & The database's data directory: everything the site stores. \\
\val{mailpit-data} & \val{mailpit} & The caught mail. \\
\val{sftp-data} & \val{sftpServer} & The reports FHHS files, when its filing is switched on. \\
\end{plaintable}

Keycloak keeps nothing: it imports the example realm into each new container, so accounts
added through its console last until the container is recreated. The script that creates
Author's \val{author} database is carried inside the file and runs on the first start of an
empty \val{pgdata} volume only.

\begin{elicitcaution}[A database from before Author needs the \val{author} database]
The database runs that script only when it initializes an \emph{empty} volume. A database
created before Author joined the stack has no \val{author} database, so \val{author-survey}
never becomes healthy and Author does not start. Run the script once, by hand; the file is
already in the container:

\begin{elicitcmd}
docker compose exec db psql -U survey -d survey \
  -f /docker-entrypoint-initdb.d/V0.0.4__CREATE_AUTHOR_DATABASE.sql
\end{elicitcmd}
\end{elicitcaution}

\subsection{What the stack leaves open}
\label{sec:compose-open}

Each of these is deliberate. It is what lets the stack start in one command and be explored
without any setup, and it is why the stack is not a site.

\begin{itemize}
  \item \textbf{Every password is documented.} The database roles, the OIDC client secrets, the
        Keycloak administrator and the SFTP account all carry the values printed in
        Chapter~\ref{sec:secrets}.
  \item \textbf{The service token is documented too.} The file gives every module the
        development value of Table~\ref{tab:secrets}, and each module warns about it at
        startup (Chapter~\ref{sec:service-token}).
  \item \textbf{Almost every service publishes a port on the host} (Table~\ref{tab:ports}):
        the database, the identity provider, the report services, the mail catcher and the
        trace collector as well as the applications. Only the SFTP server does not.
  \item \textbf{The identity provider is an example.} The bundled Keycloak runs in its
        development mode with the example realm of Chapter~\ref{sec:oidc}.
  \item \textbf{Any access code reaches a survey.} The file sets
        \opt{accessCode.autoRegister} to \val{true} on \val{survey}, so that an evaluator can
        type a code without registering a subject first.
  \item \textbf{No mail leaves the machine.} Admin sends to the mail catcher, whose web
        interface on port \val{8025} shows every invitation.
  \item \textbf{Nothing is served over \val{https}.}
\end{itemize}

\subsection{Before bringing it up}
\label{sec:compose-before}

Nothing has to change. These are the changes an evaluation most often wants:

\begin{itemize}
  \item Mount a brand, to see the site as yours (Chapter~\ref{sec:brand}). As shipped, every
        application shows the default brand its image embeds. Each of \val{survey},
        \val{admin}, \val{author} and \val{author-survey} carries the mount and its setting,
        \opt{brand.file.system.path}, commented out: uncomment both and point the volume at
        your brand directory, or at \file{elicit-brand/} in a checkout of the repository.
  \item Mount translations (Chapter~\ref{sec:i18n}). As shipped the site is in English only,
        because the images carry no other language. The same four services carry the mount
        and \opt{i18n.file.system.path}, commented out; \file{elicit-i18n/} in the repository
        is the curated set. Set \opt{i18n.offered.locales} on any application that should
        offer fewer languages than are mounted.
  \item Remove \val{fhhs} and \val{pedigree}, and Admin's dependency on \val{pedigree}, if you
        are not evaluating the Family Health History Survey. Remove \val{author} and
        \val{author-survey} if you are not evaluating authoring.
  \item Change a published host port that is already taken on your machine, and the matching
        redirect URI in the bundled realm (Chapter~\ref{sec:oidc}).
\end{itemize}

\subsection{Bringing it up}
\label{sec:compose-up}

\begin{elicitcmd}
docker compose up -d
docker compose ps
\end{elicitcmd}

The file runs the release it came with. To run another, name it in \val{ELICIT_TAG}, either
on the command line (\val{ELICIT_TAG=3.0.0-rc.1 docker compose up -d}) or once, in a
\file{.env} file beside the compose file:

\begin{elicitcmd}
ELICIT_TAG=3.0.0-rc.1
\end{elicitcmd}

A pre-release (\val{3.0.0-beta.1}, \val{3.0.0-rc.1}) is published under its exact version
only; \val{latest} and the \val{3.0} tag move only when a final version ships. A developer
sets \val{ELICIT_TAG=latest} there, which is what a local build produces. The same file is where an
evaluation sets \opt{ELICIT_SERVICE_TOKEN}, to run the stack with a service token of its own
rather than the documented one (Chapter~\ref{sec:service-token}).

That is the whole procedure. The database starts and becomes healthy, Survey creates and
migrates the shared schema, Admin migrates on top of it, and the remaining services start as
their dependencies become healthy. Go on to the first sign-in (Chapter~\ref{sec:firstsignin})
and apply a survey definition (Chapter~\ref{sec:definition}).

\begin{elicitnote}[No restart sequence is required]
Older instructions described bringing the stack up and then restarting it twice. That no
longer applies. Survey creates the schema, a survey definition is applied through the console
afterwards (Chapter~\ref{sec:definition}), and FHHS becomes ready by itself on its next
health probe. FHHS waits for Admin, not the other way round, so the console comes up on a
new database and the definition can be applied through it while FHHS is still waiting for
its survey.
\end{elicitnote}

\subsection{Resetting}

To start again from an empty database, run

\begin{elicitcmd}
./resetDatabase.sh V3
\end{elicitcmd}

which stops the stack and deletes the database's volume. Everything in the database is
lost. Without the repository, \val{docker compose down -v} does the same and deletes the caught
mail and the filed reports with it.

\subsection{From an evaluation to a real site}
\label{sec:compose-onward}

A real site is installed separately; it is not this stack with its passwords changed.

\begin{plaintable}{tab:compose-onward}{What an evaluation carries over to a real site}{@{}>{\raggedright\arraybackslash}p{0.46\linewidth} >{\raggedright\arraybackslash}p{0.48\linewidth}@{}}{\thead{Carries over} & \thead{Stays behind}}
The brand directory you made (Chapter~\ref{sec:brand}). & The compose file. \\
\addlinespace
The translations directory, and your choice of languages (Chapter~\ref{sec:i18n}). & Every
bundled service: the database, Keycloak and its realm, the mail catcher, the SFTP server and
the trace collector. \\
\addlinespace
The survey definition files you applied, once their report addresses name the site's own
services (Chapter~\ref{sec:fhhs}). & Every documented password (Chapter~\ref{sec:secrets}). \\
\addlinespace
The settings the evaluation showed your site needs (Appendix~\ref{sec:config}). & Everything
in the evaluation database: its subjects, access codes and answers. \\
\end{plaintable}

\noindent Then, in order: prepare the database (Chapter~\ref{sec:database}), configure your
identity provider (Chapter~\ref{sec:oidc}), install the applications
(Chapter~\ref{sec:nocompose}), secure the site (Chapter~\ref{sec:security}) and verify it
(Chapter~\ref{sec:verify}).

%% ================================================================== 8
\section{Installing without Compose}
\label{sec:nocompose}

This is how a real site is installed. Elicit is released as container images only,
\file{elicitsoftware/<name>} at the release's tag; no runnable Java artifact is distributed.
A site runs those images under the container platform it already has, such as Podman,
Kubernetes or a cloud container service, against its own database
(Chapter~\ref{sec:database}) and identity provider (Chapter~\ref{sec:oidc}). The bundled
compose file is for development and evaluation (Chapter~\ref{sec:compose}) and is no part of
this procedure. Start the images as below, once per module.

\begin{enumerate}
  \item Pull the module's image at the tag of the release you are installing. Every module
        runs at the same tag.
  \item Write its configuration (datasource, identity, mail, the brand path, and anything else
        your site changes) as environment variables, by the rules in
        Appendix~\ref{sec:config}, and mount the brand and translations directories read-only.
  \item Give it the site's service token, \opt{ELICIT_SERVICE_TOKEN}: one value, generated
        once, the same for every module. A site with report services also sets
        \opt{ELICIT_SERVICE_TRUSTED_ORIGINS} on Survey, Admin and FHHS
        (Chapter~\ref{sec:service-token}). A module started without the token reports itself
        not ready and says why.
  \item Run it under your container platform, publishing its container port \val{8080} to
        your reverse proxy and to nothing else (Section~\ref{sec:security-exposure}).
  \item Start \ui{Survey} first and wait for \opt{/q/health/ready} to report healthy. That is
        what confirms its migrations have completed.
  \item Start \ui{Admin} and wait for its readiness endpoint.
  \item Start \ui{FHHS}, if your site runs it, only once Admin is ready; start \ui{Pedigree}
        whenever you like (Chapter~\ref{sec:fhhs}).
  \item Secure the site before it takes a real respondent (Chapter~\ref{sec:security}).
\end{enumerate}

\begin{elicitcaution}[Nothing sequences this for you]
Unlike the compose file, your platform does not know the start order. A module started before
Survey has migrated a \emph{new} database fails its own migration on an object that does not
exist yet, and does not start. Nothing is corrupted: let Survey finish, then start the module
again.
\end{elicitcaution}

\subsection{Running more than one instance}

Instances of the same module must agree on their datasource, their brand path and their
translations: mount the same directories and set the same \opt{i18n.offered.locales} on each.
Instances of Survey that disagreed, behind one load balancer, would offer a different set of
languages depending on which one answered. Keep the proxy-forwarding settings on
(Table~\ref{tab:quarkus-http}), or sign-in redirects are built with the container's address
instead of the site's. They are safe when only the proxy reaches the instances and the proxy
replaces any forwarding header a client sent (Section~\ref{sec:security-exposure}). Survey
and Admin honor them only from the addresses in \opt{ELICIT_TRUSTED_PROXIES}: when the load
balancer is not on a private address, list it there.

%% ================================================================== 9
\section{The first sign-in}
\label{sec:firstsignin}

A new installation seeds no survey and no department. It does seed two accounts, \val{admin}
and \val{user}, so that the console can be signed into at all. (A database created before
version 3.0.0 keeps the \val{Testing Department} it was seeded with, so the department dialog
below never appears on it.)

\screenshot{04-sign-in}{Signing in. The console never asks for a password itself: it sends you
to the identity provider and takes the answer back.}

\screenshot{05-first-sign-in}{What a new installation looks like on the first sign-in: the
blocking dialog, the notice that no survey is installed, and the warning about the two seeded
accounts. All three clear during this chapter and the next.}

\begin{enumerate}
  \item Sign in to Admin with an account holding \val{elicit_admin}.
  \item A dialog blocks the console until a department exists. Follow its link to
        \ui{Departments} and create one; it is assigned to you as its creator.
  \item Rename the two seeded accounts to real people. Until both are renamed, Admin logs a
        warning at every start and shows a banner on every console screen. Under
        \opt{elicit.authorization.mode}\val{=DATABASE}, rename the matching identity-provider
        accounts to the same usernames.
  \item Write at least one message template, putting \val{<ACCESS_CODE>} wherever a subject's
        access code belongs, for example in the login link
        \val{https://<your survey host>/login/<ACCESS_CODE>}, which opens Survey with the
        code filled in.
\end{enumerate}

\screenshot{06-departments-empty}{Departments, before there are any. The dialog's link brings
you here.}

\screenshot{07-new-department}{The first department. \ui{Default message ID} is prefilled with
\val{1}, which on a new installation names a template that does not exist yet: write the
template first, or correct this field once you have.}

\screenshot{08-console-ready}{The console with a department. The dialog is gone; the two
banners stay until a survey is applied and the seeded accounts are renamed.}

\begin{elicitnote}[An account with no department can only log out]
That is deliberate. A user holding \val{elicit_user} but belonging to no department has
nothing to do in the console until an administrator assigns them one.
\end{elicitnote}

\begin{elicitcaution}[A department created in SQL is assigned to nobody]
Creating a department with a direct \val{INSERT} skips the assignment the console makes.
Assign it to its users separately, or create it through \ui{Departments}.
\end{elicitcaution}

%% ================================================================== 10
\section{Installing a survey definition}
\label{sec:definition}

A new site has no survey; surveys arrive as definition files, exported from Author and applied
in Admin.

\begin{enumerate}
  \item In Admin, open \ui{Apply Survey Definition}.
  \item Upload a \file{.elicit} file. For the Family Health History Survey that is
        \file{FHHS/family-history-survey.elicit}.
  \item The console installs the survey the first time and updates it in place afterwards,
        matching elements by their keys, so the site's respondents, subjects and reports stay
        attached to it. A file whose revision is older than the one already applied is refused:
        going back to an earlier revision is a database restore, not an upload.
  \item As part of the apply, the console asks Survey to build the reporting star schema
        of the survey just applied. Survey creates a schema of the survey's own, named
        \val{report_} followed by the survey's name in lower case (for the Family Health
        History Survey, \val{report_family_history_survey}), records the name on
        \val{survey.surveys.report_schema}, and fills it with the dimension tables, the fact
        table and the views. \val{surveyreport} itself keeps only the shared date and status
        dimensions.
  \item A new survey is in no department, and the result says so. Open \ui{Departments}, edit
        each department that runs the survey, add it under \ui{Surveys} and save. Until then
        nobody can register a subject onto it, by the form, a CSV upload or the registration
        API. A department created \emph{after} a survey is installed starts with every installed
        survey, so this step is needed only for a survey applied to departments that already
        exist, as on the first run.
\end{enumerate}

\screenshot{09-apply-definition}{Apply Survey Definition. The file identifies itself, so the
same screen installs a survey and updates one.}

\screenshot{10-apply-result}{The result. It reports what was installed and, on the last line,
that the reporting schema was rebuilt, which is why no Survey restart is needed.}

\begin{elicitnote}[No Survey restart is needed]
The rebuild happens during the apply: Admin calls Survey, at \opt{elicit.survey.url},
presenting the service token (Chapter~\ref{sec:service-token}), unless
\opt{elicit.survey.etl-build.enabled} turns the call off. The outcome is reported in the apply
result and never fails the apply. A site that restarts Survey after every apply, or runs it with
\opt{elicit.etl.enabled}\val{=false}, can turn the call off; the reporting schema is then built
when Survey next starts.
\end{elicitnote}

\begin{elicitcaution}[When the last line says the schema was not rebuilt]
The definition is installed either way. Fix the cause the line gives, then either apply the
same file again, which reconciles every record as unchanged, or restart Survey, whose start-up
build covers the same ground. The cause most often seen is two surveys at one site whose steps
or sections carry the same reporting dimension name; a refusal with \val{401} is the service
token (Chapter~\ref{sec:service-token}).
\end{elicitcaution}

To confirm the rebuild from the database:

\begin{elicitcmd}
docker compose exec db psql -U survey -d survey \
  -c "select name, report_schema from survey.surveys;" \
  -c "select count(*) from report_family_history_survey.dim_step;"
\end{elicitcmd}

\val{report_schema} must be set and the count non-zero. The fact table stays empty until
respondents finish; \val{fact_respondents} is a view over \val{survey.respondents}, so it
follows every respondent as they progress.

The schema's name belongs to the site and is never carried in a definition file. An
administrator can rename it with \ui{Rename} on Admin's \ui{Export Survey Definition} screen.
Nothing inside Elicit needs changing after a rename, but queries outside Elicit that name the
old schema do.

If your site runs FHHS, it has been reporting not-ready and answering report requests with
\val{503} until now, naming the survey it is waiting for. It goes healthy by itself on its
next probe; nothing needs restarting (Chapter~\ref{sec:fhhs}).

\elicitpart{Customizing}{Making the site its own: its brand, the languages it offers, and the
reports its surveys end with. The reports chapters are written for whoever builds a report
service, and install the two that ship with Elicit.}

%% ================================================================== 11
\section{External branding}
\label{sec:brand}

Survey and Admin each carry the Elicit default brand inside the image. To give a site its own
colors, typefaces, logos and organization name, mount a brand directory and point
\opt{brand.file.system.path} at it. Nothing is rebuilt.

\subsection{The layout}

\begin{elicitfile}[fontsize = \footnotesize]
<brand>/
  brand-config.json                name, organization, the logo file names
  brand-info.json                  description for page metadata
  theme.css                        imports the two stylesheets
  colors/brand-colors.css          --brand-* color properties
  typography/brand-typography.css  --brand-* typeface properties, @font-face rules
  typography/typography.json       the typefaces, described
  typography/fonts/                the font files, if the brand has its own
  images/                          the logo files brand-config.json names
\end{elicitfile}

Start from the repository's \file{elicit-brand/} and change only the files that differ.

\subsection{The logos}

\file{brand-config.json} names three image files, each under \file{images/}:

\begin{elicitfile}
"logos": {
  "horizontal": "HorizontalLogo.png",
  "icon": "icon-white.png"
},
"favicon": "favicon.ico"
\end{elicitfile}

\val{horizontal} is the full logo, the mark with its wording, shown on the \ui{Branding} page.
\val{icon} is the mark alone, shown in every application's header on the primary color, so it
needs a transparent ground and a color that reads on that primary. \val{favicon} is the
browser tab's icon. Any file name and any image format a browser shows will do; an SVG scales
cleanly. A name that is left out is taken from the default brand.

The icon and the favicon are the same mark at two sizes. A brand that brings its own favicon
must name its own \val{icon} too: without it the header keeps Elicit's mark beside the site's
favicon, and the \ui{Branding} page warns about it.

\subsection{Fonts}

A brand's fonts are files in the brand. \file{typography/brand-typography.css} names each one in
an \val{@font-face} rule, with an address relative to that stylesheet:

\begin{elicitfile}
@font-face {
  font-family: "IBM Plex Sans";
  font-weight: 400;
  src: url("fonts/IBMPlexSans-Regular.woff2") format("woff2");
}
\end{elicitfile}

The application serves the file from the brand, so a respondent's browser asks the site for it
and nobody else. The default brand carries IBM Plex Sans, IBM Plex Sans Arabic and IBM Plex Mono
this way, with their license beside them. A brand that names a font at another host makes every
page request it from that host, which tells that host who opened the survey; keep a brand's
fonts in the brand. A brand with no font files names a stack of system fonts instead.

\subsection{A brand may be partial}

Each file is looked for on its own: first in the mounted directory, then in a \file{brand}
directory beside the application, then in the default brand inside the image. A brand that
changes only the colors needs only \file{colors/brand-colors.css}, and a \file{theme.css} that
imports it.

A brand sets \val{--brand-*} custom properties and nothing else. Each application's own
stylesheet maps its theme's own custom properties onto them, which is what lets one brand directory serve both
applications unchanged.

\begin{elicitcaution}[A partial brand fails quietly]
A file the mount lacks, or names wrongly, is taken from the default brand without an error. The
page renders, in the default typeface or logo, and raises no error. The \ui{Branding} page
(below) is the place that shows it.
\end{elicitcaution}

\subsection{Mounting one}

Mount the directory read-only on Survey and on Admin, and set the path on each. In the compose
file:

\begin{elicitcmd}
volumes:
  - ./my-brand/:/opt/brand:ro
environment:
  brand.file.system.path: /opt/brand
\end{elicitcmd}

No application writes to the brand directory, so read-only costs nothing. The setting is per
application, so Survey and Admin may be given different brands if a site wants its
respondent-facing pages to look different from its console.

\screenshot{17-branded-header}{A branded application. The header logo, the colors and the
typefaces all come from the mounted directory.}

\subsection{Changing a brand on a running site}

Each application reads its brand once and holds it, so editing a file in the mount changes
nothing by itself. Restart the application to pick the change up. Admin's \ui{Branding} page
also has \ui{Reload brand}, which makes the console read its brand again without a restart; it
does not reach Survey, which is a separate service.

\subsection{Translated brand text}

The organization name in the header and the brand's description are translated \emph{in the
brand}, not in the applications' translation files. Add a \val{localized} block, keyed by
language tag, to \file{brand-config.json} for the name and the organization:

\begin{elicitfile}
{
  "name": "Elicit Default Brand",
  "organization": "Elicit Software",
  "localized": {
    "es-419": { "organization": "Elicit Software" },
    "ar":     { "organization": "Elicit Software" }
  }
}
\end{elicitfile}

A block of the same shape in \file{brand-info.json} translates the \val{description}. Each text
is looked up by the exact tag (\val{es-419}), then by the language alone (\val{es}), then falls
back to the untranslated value. Leave the base \val{name} as it is: the applications derive
the brand's technical key from it.

\subsection{Checking that it took}

Admin's \ui{System} section has a \ui{Branding} page. It names the brand directory that was
configured and whether it exists, and says, for each asset, whether it came from the mount, the
local directory or the default brand. It is the fastest way to find a mount that is one
directory level off.

\screenshot{13-system-branding}{The Branding page. It names the directory that resolved and,
for each asset, whether it came from the mount or from the default brand.}

%% ================================================================== 12
\section{Translations}
\label{sec:i18n}

Survey and Admin carry their own interface text, the \term{chrome}, in English inside the
image, and in no other language. Every other language is a file the site supplies, in a
translations directory mounted read-only. Elicit Software maintains a curated directory beside
each release, \file{elicit-i18n/} in the umbrella repository, which today holds Latin American
Spanish (\val{es-419}) and Arabic (\val{ar}) for Survey, Admin and the authoring tool; the
compose file mounts it.

An application offers English plus whatever is mounted for it, narrowed by
\opt{i18n.offered.locales} where that is set (Table~\ref{tab:common}). The language selector
appears in the header only when more than one language is offered, so a site that mounts
nothing never shows it.

\begin{elicitnote}[Survey content is translated separately]
This mechanism translates the chrome: buttons, labels, page titles, notifications, validation
messages and PDF headers. The survey's own content is translated too, but not from here. That
content covers:
\begin{itemize}
  \item question text and answer options;
  \item step and section names;
  \item tooltips and validation messages.
\end{itemize}
Those translations are written for one survey in the authoring tool and arrive at a site inside
its definition file. A respondent is offered a content language only when the survey publishes
it \emph{and} the site offers that language for the chrome. So the mount and
\opt{i18n.offered.locales}, which govern the chrome, also decide which content languages can
reach a respondent. Message templates are outside both mechanisms and appear in
the language they were authored in. A report service is sent the respondent's language and
writes its report in it, or not, as it chooses (Chapter~\ref{sec:reports}).
\end{elicitnote}

\subsection{The translations directory}

One directory serves every application. Each reads only its own sub-directory,
\file{<i18n.file.system.path>/<i18n.app.name>/}: by default \file{/i18n/survey/} and
\file{/i18n/admin/}, and \file{author/} for the authoring tool and its preview.

\begin{elicitfile}
elicit-i18n/                     mounted read-only at /opt/i18n
  i18n-config.json               text direction and size per language
  survey/
    translations_es_419.properties
    translations_ar.properties
  admin/
    translations_es_419.properties
    translations_ar.properties
  author/
    translations_es_419.properties
    translations_ar.properties
\end{elicitfile}

The three sub-directories hold the same languages. A language that one application has and
another lacks cannot be served consistently, so the curated set is released that way and a site
that adds a language adds it for each application it runs.

Each \file{translations_<tag>.properties} file adds its language to that application. The
applications only read the directory, on the server: nothing in it is served to a browser, and
nothing is written to it.

The curated directory holds the wording a release was built and tested with, and is versioned
with it, so mount the copy that matches the version you run. A site may mount it whole, mount
only some of its files, or mount files of its own (Section~\ref{sec:commission}). Anything
beyond the curated files is the site's to maintain.

\screenshot{18-language-selector}{The selector, open on an application offering all three
languages of the curated directory. Each row is labeled in its own language. The selector appears
only when there is something to select, so a site offering only \val{en} never sees it.}

\subsection{How a text is resolved}

Text by text: the file for the exact tag (\val{es-419}), then the file for the language alone
(\val{es}), then English. A mounted file needs only the texts that differ from English; a text
it does not carry falls back to English rather than showing blank. A file holding one corrected
text is therefore a valid way to change one wording. A text missing everywhere shows as
\val{!key!} and is logged once.

\subsection{Choosing which languages to offer}
\label{sec:language}

\begin{enumerate}
  \item Mount the translations directory read-only on Survey and Admin, and set
        \opt{i18n.file.system.path} to the mount point. The compose file carries both, commented
        out, for \file{elicit-i18n/} at \val{/opt/i18n} (Section~\ref{sec:compose-before}).
        Leave it unmounted for an English-only site.
  \item If an application should offer fewer languages than are mounted, set
        \opt{i18n.offered.locales} on it to a comma-separated list, for example
        \val{en,es-419}. Leave it unset to offer everything mounted.
  \item Restart each application. Both the directory and the list are read at startup.
\end{enumerate}

A language that is not offered is not offered anywhere: the selector hides it, a \opt{?lang=}
naming it is refused, and survey content is not served in it, because content follows the
chrome. \opt{i18n.offered.locales} cannot add a language: a tag it lists with no file behind
it is logged at startup and ignored.

\begin{elicitcaution}[The tag is punctuated with a hyphen; the file name with an underscore]
A language tag is written with a hyphen everywhere an operator sets one: \val{es-419},
\val{ar}. That covers \opt{i18n.offered.locales}, \optlit{i18n.direction.<tag>},
\optlit{i18n.font-scale.<tag>} and a \opt{?lang=} parameter. The file names in the
translations directory use an underscore instead (\file{translations_es_419.properties}). A
file named with a hyphen is not recognized as a language.
\end{elicitcaution}

The mount and \opt{i18n.offered.locales} are set per application, not per site. Survey and
Admin each read their own sub-directory and their own setting, so narrowing Admin does not
narrow Survey, and a language missing from Admin's sub-directory is not offered by the console
even when Survey offers it.

\begin{elicitnote}[Managing this from inside the applications is specified, not yet built]
Admin UC-027 specifies a \ui{Languages} screen for managing languages through the application
rather than the mounted directory and a startup property. It is not implemented in this release.
\end{elicitnote}

\subsection{Text direction and text size}
\label{sec:direction}

Arabic, Hebrew, Persian, Urdu, Pashto, Sindhi, Uyghur, Yiddish, Dhivehi and Kurdish (Sorani)
are laid out right-to-left by default. Every language is shown at the reader's own text size
unless something says otherwise. Three places can say otherwise, each winning over the one
before:
\begin{enumerate}
  \item The image's own manifest, \file{META-INF/i18n/i18n-config.json}.
  \item \file{i18n-config.json} at the root of the translations directory.
  \item A property for one tag on one application: \optlit{i18n.direction.<tag>} or
        \optlit{i18n.font-scale.<tag>} (Table~\ref{tab:common}).
\end{enumerate}
Both manifests declare \val{ar} right-to-left at a \opt{fontScale} of \val{1.15}:

\begin{elicitfile}
{ "locales": [ { "tag": "ar", "direction": "rtl", "fontScale": 1.15 } ] }
\end{elicitfile}

A site that mounts a language whose direction does not follow from its tag declares it in the
mounted \file{i18n-config.json}, or with a property when the change is for one application
only. \opt{direction} is \val{rtl} or \val{ltr}; no other direction is supported. The
applications' stylesheets mirror themselves for a right-to-left language, so no stylesheet per
language is needed.

\opt{fontScale} multiplies the reader's own base text size on pages shown in that language, and
changes nothing in any other. It exists because some scripts read smaller than Latin at the same
font size: Arabic set at the size of the Latin text beside it, in IBM Plex Sans Arabic as in
most typefaces, looks a size smaller, so a page that is comfortable in English is cramped in
Arabic. The whole page scales in proportion (text, spacing and controls together), so
nothing is clipped, and a reader who has enlarged text in their browser keeps that enlargement
on top of the scale.

To raise or lower the scale, set \optlit{i18n.font-scale.ar} or edit the mounted manifest. The
change takes effect when each application that offers the language next restarts.

\begin{elicitnote}[Every entry is optional, and a bad value is ignored]
Each value is resolved on its own:
\begin{itemize}
  \item An entry or property may set the direction, the scale, or both.
  \item Anything a layer does not declare falls back to the layer below.
  \item A language no layer mentions renders exactly as it did before any of them existed.
  \item A malformed mounted manifest is logged and the image's own stands.
  \item A \opt{fontScale} that is not a number, or that falls outside \val{0.75} to \val{2.0},
        is logged and dropped, and that language is shown unscaled. It is never clamped into
        range: a size nobody asked for is worse than none.
\end{itemize}
\end{elicitnote}

\subsection{Which language a reader gets}

In order:
\begin{enumerate}
  \item a \opt{?lang=<tag>} parameter on the address, which an invitation link can carry;
  \item the language chosen earlier in the same browser session;
  \item the browser's preferred languages, matched against those offered;
  \item English.
\end{enumerate}
An unknown tag is ignored. The choice lasts for the session only; nothing is stored against a
respondent or a user.

\subsection{A language the curated directory does not have}
\label{sec:commission}

There are two routes:
\begin{itemize}
  \item Ask Elicit Software to add the language to the curated directory. It then arrives,
        reviewed, with a later release.
  \item Supply the language yourself, and have it from the next restart.
\end{itemize}
Either way the translation starts from the same document: each application's translation
request, \file{i18n/TRANSLATION_REQUEST.md} in its repository, which contains:
\begin{itemize}
  \item the application's purpose and audience;
  \item a glossary of terms that must stay consistent or untranslated;
  \item the placeholder and markup rules;
  \item every key with its English text, where it appears, its maximum length and any flags;
  \item the return format: one \file{translations_<tag>.properties}, with the same keys in the
        same order.
\end{itemize}

To supply a language yourself, do this for each of Survey and Admin:
\begin{enumerate}
  \item Take that module's \file{TRANSLATION_REQUEST.md} from the version you run.
  \item Hand it to a translator or an AI agent.
  \item Place the returned file in that module's sub-directory of your translations directory.
  \item Restart.
\end{enumerate}
Use the request from the version you run: one from another version lacks or misnames the texts
that changed, and those then appear in English. Nobody supporting your site can see a file you
supply or a curated file you edit, so keep it with the site's configuration, and send
corrections to Elicit Software so that the curated directory carries them too.

%% ================================================================== 13
\section{Custom reports}
\label{sec:reports}

A survey can end with reports: a summary of what the respondent said, a risk estimate, a
drawing. Elicit builds one of them itself, the \term{respondent report}, which echoes back the
answers a survey author marked for it. Every other report is a \term{report service}: a small
web application your site runs, which Survey and Admin call with a respondent's id and which
answers with the report.

This chapter is the contract between Elicit and a report service: when it is called, what it
is sent, what it must answer, and where it finds the answers to report on. It is written for
whoever builds the service, the \term{report developer}. Chapter~\ref{sec:fhhs} installs the
two report services that ship with Elicit, FHHS and Pedigree, which are the working example of
everything described here.

\subsection{How a report reaches a reader}

A survey carries an ordered list of report services. Each entry is a name, a description and
an address, and it arrives at a site in the survey's definition file. Two things ask for the
reports on that list:

\begin{description}
  \item[Survey,] when a respondent finishes. The respondent lands on a report page that calls
        every service on the list, in order, and shows each answer as a card headed with the
        report's name. A \ui{Generate PDF} button on the same page sets every report in type.
  \item[Admin,] when an administrator chooses \ui{Print reports} on a finished respondent's
        row. The console calls the same services and assembles their answers into one PDF.
\end{description}

\begin{elicitdiagram}{One report request. Survey and Admin send the same request and read the
same answer; only what they do with it differs.}
  \node[elicit app, minimum width = 2.7cm] (survey) at (0, 0.9)
    {Survey\\[-1pt]{\scriptsize report page}};
  \node[elicit app, minimum width = 2.7cm] (admin) at (0, -0.9)
    {Admin\\[-1pt]{\scriptsize \ui{Print reports}}};
  \node[elicit app, minimum width = 2.7cm, minimum height = 5.2em] (svc) at (6.4, 0)
    {Report service\\[-1pt]{\scriptsize your site's}};
  \node[elicit database, minimum width = 2.7cm, minimum height = 5.2em] (db) at (12.2, 0)
    {Reporting schema\\[-1pt]{\scriptsize\val{report_<survey>}}};
  \draw[elicit flow, {Stealth[length = 5pt]}-{Stealth[length = 5pt]}] (survey.east) --
    ([yshift = 0.55cm]svc.west);
  \draw[elicit flow, {Stealth[length = 5pt]}-{Stealth[length = 5pt]}] (admin.east) --
    ([yshift = -0.55cm]svc.west);
  \node[elicit label] at (3.7, 1.45) {request \val{{"id", "language"}}};
  \node[elicit label] at (3.7, -1.45) {answer \val{{"title", "innerHTML", "pdf"}}};
  \draw[elicit flow] (svc) -- node[elicit label, above = 2pt] {reads} (db);
\end{elicitdiagram}

A survey whose list is empty has no report page: a respondent who finishes goes straight to the
survey's post-survey address.

\subsection{The request}

Each service is sent an HTTP \val{POST} to exactly the address on the list, with a JSON body:

\begin{elicitfile}
{ "id": 1042, "language": "es-419" }
\end{elicitfile}

\begin{description}
  \item[\val{id}] The respondent's id: \val{survey.respondents.id}, which is also the
        \val{respondent_id} of every row the reporting schema holds about them.
  \item[\val{language}] The tag of the language the respondent answered in, so that the
        report can be written in it. It is \val{null}, or absent, when the respondent answered
        in the survey's base language, or the survey publishes no other.
\end{description}

\begin{plaintable}{tab:report-callers}{How each caller makes the request}{@{}>{\raggedright\arraybackslash}p{0.17\linewidth} >{\raggedright\arraybackslash}p{0.37\linewidth} >{\raggedright\arraybackslash}p{0.38\linewidth}@{}}{\thead{} & \thead{Survey} & \thead{Admin}}
When & The respondent reaches the report page after finishing, and on every later sign-in
with their access code. & An administrator prints the reports of a respondent whose status
is \textbf{Finished}. \\
\addlinespace
Language & The language the respondent is reading the survey in. & The language of the
respondent's most recent answer. \\
\addlinespace
Waits & 5 seconds to connect, 60 seconds for the answer. & The REST client's defaults. \\
\addlinespace
Credentials & The site's service token, when the address is at an origin the site lists as
trusted (Section~\ref{sec:service-token-origins}). Otherwise none. & The same. \\
\addlinespace
Uses & \val{innerHTML} for the page; \val{pdf} for \ui{Generate PDF}. & \val{pdf} only. \\
\end{plaintable}

\subsection{Where a report service runs}
\label{sec:reports-network}

A report service has no user interface. It serves no pages, has no sign-in and offers nothing
for a person to open in a browser; everything a reader sees of a report is what Survey and Admin
draw from its answer. A respondent's browser never calls a report service either: Survey calls
it on the respondent's behalf, and Admin on the administrator's.

So a report service is not exposed to the Internet. Run it on a private network where its
address answers only the computers running Survey and Admin: the network the containers share
under Compose, or a firewall rule that admits those hosts and no others. A service that another
report service calls, as FHHS calls Pedigree, answers only that service.

The network is a report service's first lock. The report services that ship with Elicit have a
second: FHHS and Pedigree require the site's service token on every call
(Chapter~\ref{sec:service-token}). Survey and Admin send the token to FHHS because the site
lists FHHS's origin as trusted, and FHHS sends it to Pedigree, whose address it is configured
with.

A report service of your own is on the trusted list only if the operator puts it there
(Section~\ref{sec:service-token-origins}).

\begin{description}
  \item[Not listed,] it is called without a credential and should check none: the network
        is its only lock, and anything that can reach its address can ask it for any
        respondent's report.
  \item[Listed,] it receives the site's token in the \val{Authorization} header of
        every request, and can answer \val{401} to a request that does not carry it. But
        listing a service hands it the platform's token, with which it could itself call
        Survey, FHHS and Pedigree. A site lists a service only when it trusts that service as
        it trusts Survey.
\end{description}

Either way, do not give a report service a sign-in of its own: Survey and Admin cannot answer
one.

\begin{elicitsecurity}[Keep every report service off the Internet]
Never publish a report service's port beyond the Survey and Admin hosts, and never route it
through the reverse proxy that serves Survey to respondents, whether or not the service checks
the token. The shipped compose file publishes FHHS on the evaluation host because the shipped
Family History definition's addresses reach FHHS through that port
(Chapter~\ref{sec:fhhs}), and Pedigree for local testing. A real site does not run that file.
\end{elicitsecurity}

\subsection{The answer}

A service answers \val{200} with a JSON object of three fields:

\begin{elicitfile}[fontsize = \footnotesize]
{
  "title": "Respondent Summary",
  "innerHTML": "<table class=\"card\"><tr><th>Age</th><td>47</td></tr></table>",
  "pdf": {
    "title": "Respondent Summary",
    "landscape": false,
    "pageBreak": false,
    "content": [
      { "text": "Answers given on 12 November 2026" },
      { "table": { "headers": ["Question", "Answer"], "widths": [3, 2],
                   "body": [["Age", "47"], ["Sex at birth", "Female"]] } },
      { "svg": "<svg xmlns=\"http://www.w3.org/2000/svg\" ...>...</svg>" }
    ]
  }
}
\end{elicitfile}

\begin{description}
  \item[\val{title}] The report's own name for itself. Survey heads the card with the
        report's name from the survey's list instead, translated into the respondent's
        language, so a site renames a report in the survey, not in the service.
  \item[\val{innerHTML}] An HTML fragment that Survey places, as it is, inside the report's
        card. A fragment, not a page: no \val{<html>} or \val{<body>}. An \val{<svg>} in it
        is centered in the card.
  \item[\val{pdf}] The same report as printable content, set in type by Admin's
        \ui{Print reports} and by Survey's \ui{Generate PDF}. It carries a \val{title}, whether
        the report is \val{landscape}, whether it starts on a new page (\val{pageBreak}), and
        a list of \val{content} blocks. Each block is one of a \val{text}, an \val{svg}
        drawing, or a \val{table} with \val{headers}, relative column \val{widths} and a
        \val{body} of rows. A block carrying a drawing is set on a landscape page.
\end{description}

\begin{elicitcaution}[Without \val{pdf} the console prints nothing]
Admin's printout is made from \val{pdf} alone. A service that answers only \val{innerHTML}
shows on the respondent's report page, but is left out of the administrator's PDF without a
word.
\end{elicitcaution}

\begin{elicitcaution}[The HTML is trusted]
Survey does not sanitize \val{innerHTML}; it is part of Survey's own page. Build it from the
answers with every value escaped, and only ever list services your site controls.
\end{elicitcaution}

When a service fails, answers with an error or does not answer in time, Survey shows an error
card in that report's place and Admin prints an error block for it. The other reports are
shown and printed as usual, and the respondent's survey is finished either way.

\subsection{Where the answers are}
\label{sec:reports-data}

A report service reads the survey's \term{reporting schema}, not the transactional tables.
Survey builds one per installed survey when the definition is applied
(Chapter~\ref{sec:definition}). It holds the answers to the questions the survey's author
tagged for analysis, one row per section a respondent reached, with each analysis tag as a
column.

\subsubsection*{Find the schema on every request}

The schema's name belongs to the site: it defaults to \val{report_} followed by the survey's
name, and an administrator can rename it at any time. Look it up by the survey's key, which
is the same at every site, on every request, rather than writing the name into the service:

\begin{elicitsql}
SELECT report_schema FROM survey.surveys WHERE survey_key = :survey_key;
\end{elicitsql}

A service that does this follows a rename on its next request. Treat the result as an
identifier, quoted as one, never as text spliced into a query.

\subsubsection*{Read one respondent}

\begin{plaintable}{tab:report-views}{The objects a report reads, in the survey's reporting schema}{@{}>{\raggedright\arraybackslash}p{0.27\linewidth} >{\raggedright\arraybackslash}p{0.67\linewidth}@{}}{\thead{Object} & \thead{What a row is}}
\val{fact_sections_view} & One section a respondent reached; a repeated section has a row per
repetition. Filter on \val{respondent_id}. Fixed columns \val{step}, \val{step_instance},
\val{section}, \val{section_instance}, \val{question} and \val{item}, then one column per
analysis tag holding that respondent's answer (for the Family Health History Survey,
\val{relationship}, \val{age}, \val{gender} and the rest). \\
\addlinespace
\val{fact_respondents_view} & One respondent: \val{id} (the respondent's id), \val{status},
and the dates they were registered, first signed in and finished. \\
\addlinespace
\val{dim_step}, \val{dim_section}, \val{dim_<tag>} & The values each column can hold, for a
report that lists the options a respondent did \emph{not} choose. \\
\end{plaintable}

\begin{elicitsql}
SELECT step, section, relationship, age, gender
  FROM report_family_history_survey.fact_sections_view
 WHERE respondent_id = :id
 ORDER BY step, step_instance, section_instance;
\end{elicitsql}

The fact rows of a respondent are written when they finish, so a report is complete only for a
finished respondent, which is exactly when both callers ask for one. The schema holds answers,
not identity: a subject's name and contact details are not in it.

\subsubsection*{Connect as the reporting role}

The role script creates \val{surveyreport_user} (Chapter~\ref{sec:database}) for exactly this.
Survey grants it read access to every reporting schema it builds, and, so that it can find
them, to the two columns of \val{survey.surveys} that name each survey's schema. That grant is
this statement, which Survey's migrations run for you:

\begin{elicitsql}
GRANT SELECT (survey_key, report_schema) ON survey.surveys
   TO surveyreport_user;
\end{elicitsql}

Run it yourself, as the schema owner, if your report service connects as a role of its own
rather than \val{surveyreport_user}. It is granted column by column, so the role still reads
nothing else of the transactional schema.

FHHS connects as \val{survey_user} instead, because it also reads the subject's details for the
document it files elsewhere (Chapter~\ref{sec:fhhs}). A report that only reports needs no more
than the reporting role.

\subsection{A minimal report service}

The contract is small enough to meet in any language. This one, in Python with Flask and
psycopg, lists a respondent's answers to two tags:

\begin{elicitcmd}[fontsize = \footnotesize]
import os, html
import psycopg
from psycopg import sql
from flask import Flask, request, jsonify

app = Flask(__name__)
SURVEY_KEY = os.environ["SURVEY_KEY"]  # the survey's key, the same at every site
DB = os.environ["REPORT_DB"]  # postgresql://surveyreport_user:...@db/survey

@app.post("/report")
def report():
    respondent = int(request.get_json()["id"])
    with psycopg.connect(DB) as conn:
        schema = conn.execute(
            "SELECT report_schema FROM survey.surveys WHERE survey_key = %s",
            (SURVEY_KEY,)).fetchone()[0]
        rows = conn.execute(
            sql.SQL("SELECT relationship, age FROM {}.fact_sections_view"
                    " WHERE respondent_id = %s ORDER BY step_instance")
               .format(sql.Identifier(schema)), (respondent,)).fetchall()
    body = [[str(r or ""), str(a or "")] for r, a in rows]
    cells = "".join(f"<tr><td>{html.escape(r)}</td><td>{html.escape(a)}</td></tr>"
                    for r, a in body)
    return jsonify(
        title="Relatives",
        innerHTML=f"<table><tr><th>Relative</th><th>Age</th></tr>{cells}</table>",
        pdf={"title": "Relatives",
             "content": [{"table": {"headers": ["Relative", "Age"],
                                    "widths": [1, 1], "body": body}}]})
\end{elicitcmd}

Run it beside the applications, on the same network, and give it an address Survey and Admin
can reach and nothing outside that network can (Section~\ref{sec:reports-network}), such as
\val{http://relatives:8080/report}.

\subsection{Adding the service to a survey}

A report service is registered in the survey, never in a site's database:

\begin{enumerate}
  \item In Author, open the survey's \ui{Reports\ldots} page and choose
        \ui{Add report service\ldots}.
  \item Give it a \ui{Name}, which heads the report card, a \ui{Description}, and the
        \ui{Address} the applications will call. Translate the name and description if the
        survey publishes other languages.
  \item Drag it to its place in the list, and export the survey.
  \item Apply the exported file at the site through Admin's \ui{Apply Survey Definition}
        (Chapter~\ref{sec:definition}).
  \item Open \ui{System} $>$ \ui{Connections} in Admin. Each report service is listed by
        name and survey; check it to confirm it answers.
\end{enumerate}

\begin{elicitnote}[The list in the file is the list at the site]
Every apply replaces the survey's report list with the one in the file, and removes any service
the file no longer names. Survey calls only \val{http} and \val{https} addresses, and verifies
the certificate of an \val{https} one. An address changed by hand in the database therefore lasts only until
the next apply. Change it in Author and export again.
\end{elicitnote}

The \ui{Connections} check sends each service one read-only request and reports whether
anything answered; it never asks for a report. A service that answers that request with
\val{405}, as one accepting only \val{POST} does, is reachable.

\subsection{Doing something when a respondent finishes}
\label{sec:psa}

A report is drawn when somebody asks for it. Some sites also need something to happen at the
moment a respondent finishes: a document filed in a records system, a message sent to a
study coordinator. That is a \term{post-survey action}.

A survey carries an ordered list of them beside its reports, and like the reports they arrive
in the definition file. An author lists them on the survey's \ui{Post-survey actions\ldots} page
in Author: each a name, a description, a method, an address and whether it sends the subject's
External ID. Like the report list, the file's list is the site's list: every apply inserts,
updates and deletes actions to match it, and an action it deletes takes its delivery records with
it.

\subsubsection*{The call}

When a respondent finishes, Survey calls each action in order, waiting 5 seconds to connect and
60 for the answer:

\begin{description}
  \item[Method] \val{POST}, \val{GET} or \val{PUT}, as the author chose.
  \item[Address] As written, after its placeholders are filled and URL-encoded:
        \val{{respondentId}} with the respondent's id, and, on an action that sends the
        External ID, \val{{externalId}} with the subject's External ID.
  \item[Body] For \val{POST} and \val{PUT}, \val{{"id": <respondent id>}}. An action
        that sends the External ID adds it, and the department's code, because an External ID is
        unique only within its department:
        \val{{"id": 1042, "externalId": "MRN00123", "departmentCode": "CARD"}}. A
        \val{GET} has no body.
  \item[Answer] Any \val{2xx} is success. The body of the answer is never read: unlike a report
        service, an action returns no HTML and no \val{pdf}, and owes Survey no particular
        response.
\end{description}

\begin{elicitsecurity}[The External ID travels only over https]
An External ID usually identifies a person in another system, such as a medical record number.
An action that sends it must use an \val{https} address. Author will not save one at
\val{http}, Admin refuses a definition file carrying one, and Survey refuses to make the call,
so even a hand-edited file cannot send it in the clear.
\end{elicitsecurity}

\subsubsection*{Giving an action a credential}

Unlike a report service, a post-survey action may call a service outside the site, which may
want a credential. The credential belongs to the site, so it is set in Survey's configuration,
keyed by the action's key, and never travels in a definition file:

\begin{plaintable}{tab:psa-credentials}{Credentials for one post-survey action, on Survey}{@{}>{\raggedright\arraybackslash}p{0.47\linewidth} >{\raggedright\arraybackslash}p{0.47\linewidth}@{}}{\thead{Setting} & \thead{What it sends}}
\optlit{elicit.post-survey-action."<key>".bearer-token} & \val{Authorization: Bearer <token>}, the token as given. Supply it as an environment variable or a secret, never in a file under version control. \\
\end{plaintable}

\val{<key>} is the action's key, the second field of its record in the definition file and the
same at every site. The token is one the service issued to your site; Survey presents it as
given and never requests or refreshes one, so when the service retires a token, set the new one
and restart Survey. An action without a token of its own is sent the site's service token when
its address is at a trusted origin, which is how an action answered by one of the site's own
modules, such as FHHS, is called (Section~\ref{sec:service-token-origins}). Any other action
without a token is called without a credential. Survey's startup
diagnostics say whether each action has a credential, never its value.

\subsubsection*{When a call fails}

Survey records each delivery in \val{survey.respondent_psa}, and retries a call that failed:

\begin{plaintable}{tab:psa-states}{Delivery states Survey records}{@{}>{\raggedright\arraybackslash}p{0.20\linewidth} >{\raggedright\arraybackslash}p{0.74\linewidth}@{}}{\thead{State} & \thead{Meaning}}
\val{PENDING}, \val{RESENDING} & Being called, for the first time or again. \\
\val{DELIVERED} & The service answered \val{2xx}. \\
\val{CALL_FAILED} & No answer in time, no connection, or another status; retried. \\
\val{CALL_REFUSED} & Refused before the call: an External-ID action at an \val{http} address, an
address that needs an External ID the subject does not have, or a placeholder it may not use.
Never retried, because a retry cannot change the outcome. \\
\end{plaintable}

Every \opt{elicit.post-survey-action.retry.interval} (15 minutes unless changed), Survey calls
each \val{CALL_FAILED} delivery again, until it is delivered or has had
\opt{elicit.post-survey-action.retry.max-attempts} attempts (24, the first included). Each
retried delivery is claimed first, so several Survey instances never call it twice, and
\opt{elicit.post-survey-action.retry.enabled} turns the schedule off (Table~\ref{tab:survey}).

\begin{elicitcaution}[Survey retries the call, not the service's work]
Once a call is answered, what the service does next is its own. If it records its own progress
on the same delivery, as FHHS does with \val{STARTED}, \val{COMPLETED} and \val{FAILED}, Survey
never overwrites that and never retries it. A service that must finish reliably answers at once,
does the work in the background, and keeps its own record and its own retries, as FHHS does for
the documents it files (Section~\ref{sec:fhhs-sftp}).
\end{elicitcaution}

%% ================================================================== 14
\section{The Family History reports}
\label{sec:fhhs}

The Family Health History Survey ends with three reports, all answered by one report service,
\ui{FHHS}, which asks a second, \ui{Pedigree}, to draw the family tree. Together they are the
reference implementation of the contract in Chapter~\ref{sec:reports}: a site that runs the
survey installs them from this chapter, and a report developer reads them as a working
example.

Both are optional. A site that does not run the Family Health History Survey runs neither.

\begin{elicitdiagram}{The Family History reports. FHHS answers three report services and one
post-survey action; Pedigree only ever hears from FHHS.}
  \node[elicit app, minimum width = 2.4cm] (survey) at (0, 1.0) {Survey};
  \node[elicit app, minimum width = 2.4cm] (admin) at (0, -1.0) {Admin};
  \node[elicit app, align = left, minimum width = 4.2cm, inner ysep = 7pt] (fhhs) at (5.6, 0)
    {\textbf{FHHS}\\[3pt]
     \texttt{\scriptsize /proband/report}\\[-1pt]
     \texttt{\scriptsize /casummary/report}\\[-1pt]
     \texttt{\scriptsize /pedigree/report}\\[-1pt]
     \texttt{\scriptsize /familyhistory/generate}};
  \node[elicit app, minimum width = 2.6cm] (ped) at (11.4, 1.0)
    {Pedigree\\[-1pt]\texttt{\scriptsize /svg}};
  \node[elicit service, minimum width = 2.6cm] (sftp) at (11.4, -1.0)
    {SFTP server\\[-1pt]{\scriptsize optional}};
  \node[elicit database, minimum width = 4.2cm] (db) at (5.6, -2.7)
    {\val{report_family_history_survey}};
  \draw[elicit flow] (survey.east) -- node[elicit label, above = 2pt, sloped] {reports, action}
    ([yshift = 0.5cm]fhhs.west);
  \draw[elicit flow] (admin.east) -- node[elicit label, below = 2pt, sloped] {reports}
    ([yshift = -0.5cm]fhhs.west);
  \draw[elicit flow] ([yshift = 0.5cm]fhhs.east) -- node[elicit label, above = 2pt, sloped]
    {pedigree table} (ped.west);
  \draw[elicit flow] ([yshift = -0.5cm]fhhs.east) -- node[elicit label, below = 2pt, sloped]
    {PDF + index} (sftp.west);
  \draw[elicit flow] (fhhs) -- node[elicit label, right = 2pt] {reads} (db);
\end{elicitdiagram}

\subsection{What they answer}

The Family History definition names these FHHS addresses; Survey and Admin call them, never a
browser.

\begin{plaintable}{tab:fhhs-endpoints}{FHHS's addresses, as the survey's definition lists them}{@{}>{\raggedright\arraybackslash}p{0.30\linewidth} >{\raggedright\arraybackslash}p{0.17\linewidth} >{\raggedright\arraybackslash}p{0.45\linewidth}@{}}{\thead{Address} & \thead{Listed as} & \thead{What it answers}}
\val{/proband/report} & Report \emph{Respondent} & The respondent's own history, as a table. \\
\addlinespace
\val{/casummary/report} & Report \emph{Cancer Summary} & Each relative's cancers and ages at
diagnosis. \\
\addlinespace
\val{/pedigree/report} & Report \emph{Family Pedigree} & The family tree, drawn by Pedigree,
as an SVG on a landscape page of its own. \\
\addlinespace
\val{/familyhistory/generate} & Post-survey action & Files all three as one PDF on an SFTP
server, when the site turns that on (Section~\ref{sec:fhhs-sftp}). \\
\end{plaintable}

Each report reads the survey's \val{fact_sections_view} for the respondent, resolving the
schema's name from \val{survey.surveys.report_schema} on every request, exactly as
Section~\ref{sec:reports-data} describes. The source is in the FHHS repository, one package
per report (\file{proband}, \file{casummary}, \file{pedigree}); it is the place to start a
report of your own.

\subsection{The two services in the compose stack}

The shipped compose file already runs both, as the services \val{fhhs} and \val{pedigree}.
FHHS waits for Admin to be ready before it starts (Chapter~\ref{sec:topology} explains why), and
connects with the same database passwords as Survey (Chapter~\ref{sec:secrets}). The file
gives both services the development service token, and lists FHHS's published address as a
trusted origin on \val{survey}, \val{admin} and \val{fhhs} (Chapter~\ref{sec:service-token}).
Bring the stack up as in Chapter~\ref{sec:compose}, then:

\begin{enumerate}
  \item Watch FHHS report \emph{not ready}. Until the survey is installed it logs one warning
        naming the survey it is waiting for, its readiness check gives the same reason, and
        it refuses report requests with \val{503} and that message.
  \item Make sure the survey's report addresses reach FHHS from the Survey and Admin containers
        (see \emph{The report addresses} below).
  \item Apply \file{FHHS/family-history-survey.elicit} in Admin
        (Chapter~\ref{sec:definition}).
  \item FHHS becomes ready on its next readiness probe. Nothing needs restarting.
  \item Add the survey to your department on \ui{Departments} (Chapter~\ref{sec:definition}),
        so subjects can be registered onto it.
  \item In Admin, open \ui{System} $>$ \ui{Connections} and check the three report services and
        the post-survey action.
  \item Walk the survey to the end as a respondent. The report page shows three cards, the
        last of them the pedigree.
\end{enumerate}

\subsection{The two services on a site without Compose}

On a site installed as Chapter~\ref{sec:nocompose} describes, start the two services yourself:

\begin{enumerate}
  \item Start \file{elicitsoftware/pedigree} at the release's tag. It listens on \val{8080},
        keeps no data and reads one setting, the site's service token
        (Chapter~\ref{sec:service-token}). Its health endpoint is \opt{/healthz}.
\end{enumerate}

\begin{elicitcmd}[fontsize = \footnotesize]
docker run -d --name pedigree \
  -e ELICIT_SERVICE_TOKEN='<service token>' \
  elicitsoftware/pedigree:<version>
\end{elicitcmd}

\begin{enumerate}[resume]
  \item Start \file{elicitsoftware/fhhs} at the same tag, \emph{after Admin is ready}, with the
        datasource settings Survey uses and the settings in Table~\ref{tab:fhhs}. At minimum:
\end{enumerate}

\begin{elicitcmd}[fontsize = \scriptsize]
docker run -d --name fhhs \
  -e QUARKUS_DATASOURCE_JDBC_URL='jdbc:postgresql://db.example.org:5432/survey' \
  -e QUARKUS_DATASOURCE_OWNER_JDBC_URL='jdbc:postgresql://db.example.org:5432/survey' \
  -e ELICIT_DB_PASSWORD='<survey_user password>' \
  -e ELICIT_OWNER_DB_PASSWORD='<elicit_owner password>' \
  -e ELICIT_SERVICE_TOKEN='<service token>' \
  -e ELICIT_SERVICE_TRUSTED_ORIGINS='http://fhhs.internal:8080' \
  -e PEDIGREE_URL='http://pedigree.internal:8080/svg' \
  elicitsoftware/fhhs:<version>
\end{elicitcmd}

The service token is the value Survey and Admin hold. The trusted origin is FHHS's own
address, as the survey's report addresses give it, because FHHS asks those addresses for the
reports when it builds the document it files. Set the same origin on Survey and Admin, or they
call FHHS without the token and FHHS refuses them
(Section~\ref{sec:service-token-origins}).

Keep the connection options of the shipped JDBC address (Appendix~\ref{sec:config}) when you
point FHHS at your own server. FHHS reports ready at \opt{/q/health/ready} once the survey is
applied, as above.

\subsection{The report addresses}

The survey's definition file lists FHHS's three reports and its action at
\val{http://host.docker.internal:8082/…}: FHHS's published port on the evaluation host, which
the Survey and Admin containers reach through Docker Desktop's host alias. That works for the
shipped compose file on a developer's machine and nowhere else.

A real site gives the four addresses the name the applications know FHHS by on their own
network, for example \val{http://fhhs.internal:8080/proband/report}. Change them in Author, on the
survey's \ui{Reports\ldots} page for the three reports and on its \ui{Post-survey actions\ldots}
page for the action, and export the survey, then apply that file instead of the shipped one.
A change made in the database is undone by the next apply.

The origin of those four addresses, \val{http://fhhs.internal:8080} in the example, is the one
to list in \opt{ELICIT_SERVICE_TRUSTED_ORIGINS} on Survey, Admin and FHHS. The two must agree:
an address at an origin that is not listed is called without the service token, and FHHS
answers \val{401} (Section~\ref{sec:service-token-origins}).

\subsection{FHHS settings}

FHHS takes the datasource and observability settings Survey does, with the same names and
passwords (Appendix~\ref{sec:config}). Its own are these.

\begin{settings}{tab:fhhs}{FHHS settings}
\opt{pedigree.url} & --- & \chgset & Pedigree's drawing endpoint, for example
\val{http://pedigree:8080/svg}. No value is shipped; the compose file sets it.
Environment: \opt{PEDIGREE_URL}. \\
\opt{elicit.service.token} & --- & \chgset & The site's service token, which FHHS requires on
every path outside \opt{/q} and presents to Pedigree and to its own report addresses
(Chapter~\ref{sec:service-token}). At least 32 characters; no value is shipped. A secret.
Environment: \opt{ELICIT_SERVICE_TOKEN}. \\
\opt{elicit.service.trusted-origins} & --- & \chgset & The origin of FHHS's own report
addresses, as the survey's definition gives them, for example
\val{http://fhhs.internal:8080}. A comma-separated list. Unset, FHHS sends the token to
Pedigree only, and its own report addresses refuse it when it builds the document it files
(Section~\ref{sec:service-token-origins}). Environment:
\opt{ELICIT_SERVICE_TRUSTED_ORIGINS}. \\
\opt{family.history.survey.key} & \val{5e91c606-59a1-450a-a8d7-2f1530ff472b} & \chgkeep & The
key of the survey FHHS serves. It is the key in the shipped definition file, the same at every
site. \\
\opt{family.history.upload.psa.key} & \val{9f57b644-fecd-557d-a466-07ec19bb95c4} & \chgkeep &
The key of the post-survey action FHHS records its uploads against. \\
\opt{family.history.upload.psa.id} & --- & \chgmay & The same action by its id at this site.
Set, it overrides the key; leave it unset. \\
\opt{family.history.async.threads} & \val{2} & \chgmay & How many family-history documents
FHHS builds at once. \\
\opt{app.base.url} & \val{http://localhost:8080} & \chgmay & The survey's public address,
printed at the foot of the filed PDF. Set it when filing is on. Environment:
\opt{APP_BASE_URL}. \\
\opt{family.history.sftp.enabled} & \val{false} & \chgmay & Whether the post-survey action
files documents on an SFTP server (Section~\ref{sec:fhhs-sftp}). Left \val{false}, the
action answers and files nothing. \\
\opt{family.history.sftp.host} & --- & \chgset & The SFTP server, when filing is on. \\
\opt{family.history.sftp.port} & \val{22} & \chgmay & \\
\opt{family.history.sftp.path} & --- & \chgset & The directory documents are written to. \\
\opt{family.history.sftp.username} & --- & \chgset & The account FHHS signs in as. \\
\opt{family.history.sftp.password} & --- & \chgmay & Its password, when the server accepts
one. A secret. \\
\opt{family.history.sftp.privateKey} & --- & \chgmay & A private key instead: a path to a
mounted key file, or the key itself. Takes precedence over the password. A secret. \\
\opt{family.history.sftp.knownHosts} & --- & \chgset & The server's host key, as
\val{ssh-keyscan} prints it. FHHS refuses a server whose key it cannot verify. \\
\opt{family.history.sftp.strictHostKeyChecking} & \val{true} & \chgkeep & Whether to insist on
that verification. The compose file turns it off for the bundled server, which makes a new key
on every start; never turn it off for a real one. \\
\opt{family.history.sftp.timeout} & \val{30000} & \chgmay & How long, in milliseconds, to wait
for the server. \\
\opt{family.history.sftp.xml.template} & an index document & \chgmay & The index written
beside each PDF (Section~\ref{sec:fhhs-sftp}). \\
\end{settings}

\subsection{Filing the family history elsewhere}
\label{sec:fhhs-sftp}

When \opt{family.history.sftp.enabled} is \val{true}, FHHS answers the survey's post-survey
action by building one PDF of every report on the survey's list and uploading it to the SFTP
server, beside an index document a records system can file it by:

\begin{elicitfile}
<path>/<xid>.pdf
<path>/<xid>-index.xml
\end{elicitfile}

The index comes from \opt{family.history.sftp.xml.template}, in which these placeholders are
replaced from the respondent's record: \val{{RespondentId}}, \val{{Xid}},
\val{{SurveyId}}, \val{{FirstName}}, \val{{MiddleName}}, \val{{LastName}},
\val{{Dob}}, \val{{Email}}, \val{{Phone}}, \val{{DepartmentName}},
\val{{DepartmentID}}, \val{{AccessCode}}, \val{{Status}}, \val{{Created}} and
\val{{Finalized}}. The compose file carries a commented example.

FHHS signs in to the server with a password or with a private key (Table~\ref{tab:fhhs}); prefer
the key. Generate a key pair for the site, give the server its public half, and supply the
private half to FHHS as a secret (Section~\ref{sec:security-secrets}). The pair under
\file{sftp_ssh/} in the repository belongs to the bundled SFTP service and is published with
it: no server of a site's own accepts it.

FHHS answers the action at once and does the work in the background. It records every upload
against the respondent in \val{survey.respondent_psa}, as started, completed or failed, and
every fifteen minutes retries the failed ones, up to fifty attempts. That is the pattern
Section~\ref{sec:psa} recommends for any action that must arrive. The two retries never
overlap: Survey calls FHHS again only when its call never reached FHHS (\val{CALL_FAILED});
once FHHS has answered, the filing and its retries are FHHS's.

\subsection{Pedigree}

Pedigree is an R service built on the \emph{Pedixplorer} package. The endpoint that matters is
\val{POST /svg}. It takes a tab-separated pedigree table, one row per family member, as the
multipart field \val{ped}, and answers with the drawing as an SVG. The table's columns are
\val{famid}, \val{id}, \val{sex}, \val{dadid} and \val{momid}, and optionally \val{deceased},
\val{proband}, \val{affection}, \val{avail}, \val{display_id} and \val{cancer_label}. FHHS
builds the table from the respondent's answers and waits 15 seconds for the drawing. A request
body over 1~MB is refused with \val{413}.
\val{GET /healthz} answers \val{{"status":"ok"}} when the service is up.

Every request but \opt{/healthz} must carry the site's service token, which Pedigree reads
from the environment variable \opt{ELICIT_SERVICE_TOKEN}; it has no other setting. Started
without the variable, Pedigree answers \val{503} to every drawing request
(Chapter~\ref{sec:service-token}).

Pedigree is a report service's helper, not a report service: nothing in Elicit calls it but
FHHS, so it never appears on a survey's report list.

\elicitpart{Running the site}{Securing the site, checking that every part of the installation
works, what to do when one does not, and where to read further.}

%% ================================================================== 15
\section{Securing the site}
\label{sec:security}

An Elicit site holds what respondents say about themselves, and often about their health. The
applications protect it in the ways this chapter describes, and leave four things to the
deployment: a private network for everything that is not a browser's business, a reverse proxy
that terminates TLS, an identity provider, and somewhere to keep secrets. This chapter says
what each control is and why it is there, and closes with a checklist
(Section~\ref{sec:security-checklist}) to work through before the site takes its first real
respondent.

\begin{elicitsecurity}[The bundled stack is not a production site]
The compose file, and the database, identity provider, mail catcher, SFTP server and trace
collector bundled with it, exist so that a developer or an evaluator has a working site in one
command. They publish every port, ship documented passwords, and prefer ease of use to
strictness throughout. The bundled Keycloak realm is an example of an OIDC provider and nothing
more. Nothing in this chapter hardens them. It describes a site installed against its own
database, identity provider and network (Chapter~\ref{sec:nocompose}).
\end{elicitsecurity}

\subsection{What faces the Internet}
\label{sec:security-exposure}

Two things on a site are for people: Survey, for respondents, and Admin, for administrators.
Everything else is for the applications.

\begin{plaintable}{tab:exposure}{Who reaches each service}{@{}>{\raggedright\arraybackslash}p{0.22\linewidth} >{\raggedright\arraybackslash}p{0.30\linewidth} >{\raggedright\arraybackslash}p{0.40\linewidth}@{}}{\thead{Service} & \thead{Reached by} & \thead{So it is}}
Survey & Respondents, in a browser. Admin, for the reporting schema. & Published through the
reverse proxy, with the paths below refused. \\
\addlinespace
Admin & Administrators, in a browser. A machine client, if the site imports subjects through
the integration API. & Published through the reverse proxy. Where every administrator works
inside the organization, published to its network only. \\
\addlinespace
A report service, FHHS included & Survey and Admin. & On the private network, answering those
hosts only (Section~\ref{sec:reports-network}). \\
\addlinespace
Pedigree & FHHS. & On the private network, answering FHHS only. \\
\addlinespace
PostgreSQL & The applications, and the people who administer it. & On the private network. \\
\addlinespace
The identity provider & Administrators' browsers, and Admin. & Whatever your organization
already does with it. \\
\addlinespace
The SMTP relay & Admin. & On the private network. \\
\addlinespace
The trace collector, and a metrics scraper & The applications; operators. & On the private
network. \\
\end{plaintable}

\subsubsection*{What the proxy refuses}

Some paths on Survey and Admin are for the site's own applications and operators, not for the
people the proxy serves. The reverse proxy refuses them, whatever the application behind it
would have answered:

\begin{description}
  \item[On Survey,] every path under \opt{/api} except \opt{/api/pdf/download}, which a
        respondent's browser uses for \ui{Generate PDF}. The rest is how Admin asks Survey to
        build, rename and drop reporting schemas and to print the respondent report. Admin
        reaches Survey at \opt{elicit.survey.url}, on the private network, not through the
        proxy.
  \item[On Survey and Admin,] every path under \opt{/q}: health, readiness and metrics
        (Section~\ref{sec:security-observability}).
\end{description}

Survey also requires the site's service token on those \opt{/api} paths
(Chapter~\ref{sec:service-token}). The proxy rule stays: a path that is not for the people the
proxy serves should not reach them at all, and two locks are better than one.

\subsubsection*{Only the proxy reaches the applications}

Survey and Admin trust the forwarding headers a proxy sets, so that sign-in redirects and links
carry the site's public address (Table~\ref{tab:quarkus-http}). That trust is safe only when
the proxy is the one sending them. Let nothing but the proxy reach an application's port, and
have the proxy replace any \val{Forwarded} and \val{X-Forwarded-*} header
a client sent rather than pass it on.

Survey and Admin also check where the headers come from: each honors them only from the
addresses in \opt{ELICIT_TRUSTED_PROXIES}, which by default are loopback and the private
ranges. Set it on both to your reverse proxy's address or network, so that another machine on
the private network cannot choose the address Admin builds its sign-in redirects with, or the
client address either application logs.

\subsubsection*{Limit the rate of requests}

Survey is open to anyone who has its address, and an access code is all that stands between
that address and a survey (Section~\ref{sec:security-access-codes}). Limit the number of
requests one client address may make at the proxy, on Survey in particular, so that codes
cannot be guessed at speed.

\subsection{Transport security}
\label{sec:security-transport}

Each application speaks plain HTTP on port \val{8080} inside its container and expects the
reverse proxy to terminate TLS.

\begin{itemize}
  \item \textbf{Browsers use \val{https} only.} Redirect \val{http} to \val{https} at the
        proxy, for Survey and for Admin. The invitation link in every message template is an
        \val{https} address (Chapter~\ref{sec:firstsignin}).
  \item \textbf{The identity provider's issuer address is \val{https}}
        (\opt{quarkus.oidc.auth-server-url}).
  \item \textbf{Outbound certificates are verified.} Leave \opt{quarkus.tls.trust-all} at
        \val{false}. A service signed by a private authority is trusted by adding the authority,
        never by turning verification off (Table~\ref{tab:quarkus-http} and the note below it).
  \item \textbf{Mail is sent over TLS.} Set \opt{quarkus.mailer.tls} to \val{true} unless the
        relay is on the same private network. Invitations carry access codes.
  \item \textbf{The database connection is encrypted} when the database is on another host:
        add \val{sslmode=verify-full} to the parameters of both connection URLs
        (Table~\ref{tab:quarkus-db}).
  \item \textbf{A report service may be called over \val{http} only inside the private
        network.} A post-survey action that leaves the site uses \val{https}, and one that
        sends the External ID is refused at anything else (Section~\ref{sec:psa}).
  \item \textbf{The SFTP server's host key is verified}
        (Section~\ref{sec:security-reports}).
\end{itemize}

\subsubsection*{Response headers}

Survey and Admin send these with every response. The proxy passes them through; it need not
add them, and must not strip them.

\begin{plaintable}{tab:security-headers}{Response headers Survey and Admin send}{@{}>{\raggedright\arraybackslash}p{0.44\linewidth} >{\raggedright\arraybackslash}p{0.50\linewidth}@{}}{\thead{Header} & \thead{Why}}
\val{Strict-Transport-Security: max-age=31536000; includeSubDomains} & A browser
that has seen the site once never asks for it, or for a name beneath it, over \val{http}
again. \\
\addlinespace
\val{X-Content-Type-Options: nosniff} & A browser does not guess the type of a file
it is sent. \\
\addlinespace
\val{Content-Security-Policy: frame-ancestors} & Another site cannot show yours
inside a frame. Survey allows its own pages to (\val{'self'}); Admin allows none
(\val{'none'}, with \val{X-Frame-Options: DENY}). \\
\addlinespace
\val{Referrer-Policy} & The address of a page is never sent to another site:
\val{no-referrer} on Survey, \val{same-origin} on Admin. \\
\end{plaintable}

\noindent Each is a setting of the form
\optlit{quarkus.http.header."<name>".value}, which a site overrides like any other
(Appendix~\ref{sec:config}). The one most likely to need it is
\val{Strict-Transport-Security}: a site whose host name has sibling names still
served over plain \val{http} drops \val{includeSubDomains}. A fuller
\val{Content-Security-Policy} is worth having and is specific to the site: it must
allow whatever the site's brand loads. Add it at the proxy in report-only mode against the
mounted brand before enforcing it.

\subsection{The identity provider and its accounts}
\label{sec:security-identity}

Sign-in to the console is decided by the identity provider, so the console is as well
protected as the provider is.

\begin{itemize}
  \item \textbf{Use your organization's provider.} The bundled realm is an example
        (Chapter~\ref{sec:oidc}); a site never imports it.
  \item \textbf{Make the redirect address exact:} the console's public \val{https} address and
        nothing wider.
  \item \textbf{Allow the console's client the authorization code flow only.} People sign in
        on the provider's own page, so the client needs no other way of obtaining a token: turn
        off the grant that exchanges a user name and password directly, and every other flow
        the provider offers on the client.
  \item \textbf{Require a second factor} for every account that holds \val{elicit_admin}, by
        whatever means the provider offers.
  \item \textbf{Limit failed sign-ins and keep a record.} Have the provider lock an account
        out after repeated failures, and record sign-ins and administrative changes where the
        people who run the site can read them.
  \item \textbf{Grant roles narrowly} (Table~\ref{tab:roles-oidc}). Most people who use the
        console need \val{elicit_user} and the departments they work in, and see only those
        departments' subjects.
  \item \textbf{Keep \val{elicit_importer} for machines.} A machine client that imports
        subjects is a client of its own in the provider, with its own secret and that one
        role. No person holds it, and the console no longer grants it to one: \val{elicit_admin}
        and \val{elicit_user} do not imply it, and the registration API accepts no other role.
        Enter the client in the console as well, under \ui{Admin} $>$ \ui{Import Clients},
        with the name its tokens carry (for a Keycloak service account,
        \val{service-account-} followed by the client ID), and assign it the departments it
        registers into. The registration API answers 403 to a client with the role and no
        active import client record, and to a token that carries \val{elicit_admin} or
        \val{elicit_user} beside it.
  \item \textbf{Rename the seeded accounts} \val{admin} and \val{user}
        (Chapter~\ref{sec:firstsignin}).
  \item \textbf{Remove a person at the provider first.} When someone leaves, disable their
        account in the identity provider, then remove their departments in \ui{Users}.
  \item \textbf{Prefer roles from the provider.} Under
        \opt{elicit.authorization.mode}\val{=DATABASE} the console also reads roles from its own
        table, so whoever can write to the database can grant them.
\end{itemize}

\subsection{Respondents and access codes}
\label{sec:security-access-codes}

A respondent has no account. The access code is their whole credential: whoever holds it can
open that subject's survey and, once it is finished, its reports.

\begin{itemize}
  \item \textbf{Decide on \opt{accessCode.autoRegister}.} Left at \val{false}, only a code
        the console issued reaches a survey. Set to \val{true}, any code of 5 to 12 characters
        typed in creates a respondent and starts the survey: anyone who can reach the login page
        can take it, which is what a survey open to the public wants. Such a respondent was never
        registered: Admin lists it as \ui{Auto-registered}, with no name, contact details or
        department, to every console account; a post-survey action receives no External ID for
        it; and the code the respondent chose is the only thing that brings them back to their
        answers. The shipped compose file turns it on.
  \item \textbf{Send a code to its subject and nobody else.} Admin's messages are the intended
        route. A code read out, printed or forwarded is a code given away.
  \item \textbf{The invitation link carries the code in its address.} Survey keeps the link
        out of its own access log and replaces the address with \opt{/login} once the page has
        opened, so the code does not stay in the browser's history. The proxy in front of it
        still sees the request as it arrived: read Section~\ref{sec:security-observability}
        before deciding what the proxy logs and who may read it.
\end{itemize}

\begin{elicitcaution}[Auto-registration can skew a survey's results]
With \opt{accessCode.autoRegister} on in production, nothing ties a respondent to a person. The
same person can take the survey again and again, each time with a new access code, and an
automated bot can type codes and enter answers as fast as it likes. Either one skews the
survey's results, and the console cannot tell such respondents apart from genuine ones. Turn it
on only for a survey whose results can bear that, and leave it \val{false} where every
respondent must be a registered subject.
\end{elicitcaution}

\subsection{The database}
\label{sec:security-database}

Everything a respondent says is in the database, with the names, dates of birth and contact
details Admin holds for each subject.

\begin{itemize}
  \item \textbf{No address outside the application hosts.} The database answers Survey, Admin,
        FHHS and the people who administer it.
  \item \textbf{Each application connects as its own role} (Table~\ref{tab:roles}). A report
        service of the site's own connects as \val{surveyreport_user}, which reads the reporting
        schemas, and never as the schema owner.
  \item \textbf{The schema owner's password is the most sensitive secret of the site.} Every
        module holds it, because every module runs its own migrations, and whoever has it can
        read and change everything Elicit stores.
  \item \textbf{No Elicit role is a superuser.} The schema owner owns the Elicit schemas and
        can create a schema in the Elicit database; it can do nothing else in the cluster
        (Chapter~\ref{sec:database}). Check it on a cluster an earlier release prepared:
        \val{rolsuper} is false for \val{elicit_owner} and \val{dbowner} in \val{pg_roles}.
  \item \textbf{Backups are protected as the database is:} encrypted, and readable by the same
        few people.
\end{itemize}

\subsection{Secrets}
\label{sec:security-secrets}

Chapter~\ref{sec:secrets} lists the credentials that ship with a documented value. A site also
has secrets of its own: the service token its modules present to each other
(Chapter~\ref{sec:service-token}), the mail relay's password, each post-survey action's bearer
token
(Section~\ref{sec:psa}), and FHHS's SFTP password or private key.

\begin{itemize}
  \item \textbf{Supply every secret as configuration,} through an environment variable or
        your platform's secret store (Appendix~\ref{sec:config}). Never write one into a file
        under version control, and never build one into an image.
  \item \textbf{A secret is changed in two places:} where it is issued, and wherever it is
        presented. Change both, then restart the modules that present it.
  \item \textbf{Change a secret when someone who knew it leaves,} and after any incident.
  \item \textbf{The applications never show one.} Survey's startup diagnostics and Admin's
        \ui{System} pages report a secret as present or absent, and mask tokens.
\end{itemize}

\subsection{What the site loads}
\label{sec:security-content}

Three things reach a site from outside it and are then shown to people or acted on.

\begin{description}
  \item[A survey definition] carries the text and markup respondents will read, and the
        addresses Survey and Admin will call for the survey's reports and post-survey actions.
        Apply only a file exported from Author by an author the site trusts. Read the
        addresses each new file brings before applying it; afterwards Admin's
        \ui{Connections} page shows the report services among the site's outbound
        dependencies (Chapter~\ref{sec:verify}).
  \item[The brand directory] supplies the style sheet, images and names on every page. Mount
        it read-only, as shipped, and let only operators write to it
        (Chapter~\ref{sec:brand}). A brand is code as much as decoration: its style sheet can
        restyle or hide anything on a page, so review a brand as you would a release. Have it
        carry its own font files rather than import them from another site, which would send
        every reader's address there.
  \item[The translations directory] supplies the text of every page in each language it
        holds. The same rule applies (Chapter~\ref{sec:i18n}).
\end{description}

\subsection{Report services and post-survey actions}
\label{sec:security-reports}

\begin{itemize}
  \item \textbf{A report service answers Survey and Admin only}
        (Section~\ref{sec:reports-network}). On a site, the Family History report addresses
        name FHHS on the private network (Chapter~\ref{sec:fhhs}).
  \item \textbf{FHHS and Pedigree require the site's service token,} and Survey, Admin and
        FHHS send it only to the origins the site lists as trusted. List no service the site
        does not run and trust (Section~\ref{sec:service-token-origins}).
  \item \textbf{A post-survey action outside the site has a credential of its own:} a bearer
        token set on Survey, never carried in the definition file (Section~\ref{sec:psa}).
  \item \textbf{FHHS files documents over SFTP with the server's host key verified.} Set
        \opt{family.history.sftp.knownHosts}, leave
        \opt{family.history.sftp.strictHostKeyChecking} at \val{true}, and prefer a private key
        to a password (Table~\ref{tab:fhhs}). The filed documents and their index hold
        personal data, the access code among it if the index template includes it
        (Section~\ref{sec:fhhs-sftp}): the server and the directory are as sensitive as the
        database.
\end{itemize}

\subsection{Health, metrics, logs and traces}
\label{sec:security-observability}

\begin{description}
  \item[Health and metrics] under \opt{/q} answer without a sign-in, so that a container
        platform can probe them and a scraper can read them. That is why the proxy refuses
        \opt{/q}: read them on the private network only.
  \item[Access logs] are on as shipped (\opt{quarkus.http.access-log.enabled}). Survey's
        records the method and path of each request, without the query string or the
        referring page, and leaves out requests for an invitation link altogether, so it holds
        no access codes. The proxy's own log does, unless you configure it the same way: leave
        the path of \opt{/login/} requests out of it, or restrict who reads it and keep it only
        as long as your policy needs.
  \item[Application logs] can name respondents and subjects. Leave the log levels as shipped
        on a site (Table~\ref{tab:quarkus-obs}), and treat the logs, and wherever they are
        shipped to, as holding personal data.
  \item[Traces] go to the collector at \opt{quarkus.otel.exporter.otlp.endpoint} and carry
        request addresses too. Send them to a collector your site protects. The bundled
        Jaeger has no sign-in and is for evaluation.
\end{description}

\subsection{Releases}
\label{sec:security-releases}

\begin{itemize}
  \item \textbf{Run a named release,} the same one for every module, never \val{latest}
        (Chapter~\ref{sec:nocompose}).
  \item \textbf{Move to a new release when one is published,} by the procedure for that
        release in \file{DeploymentScript.md}. Fixes arrive in releases; earlier lines are not
        patched.
  \item \textbf{Keep what Elicit runs on current:} PostgreSQL, the container runtime, the
        proxy and the host.
\end{itemize}

\subsection{The checklist}
\label{sec:security-checklist}

Work through this before the site takes its first real respondent, and again after every
upgrade or change of network. Each line names the section that explains it. A line about
something the site does not run does not apply.

\subsubsection*{Exposure}
\begin{itemize}[label = \elicitcheck]
  \item The site is not the bundled compose stack, and uses none of its bundled services
        (Chapter~\ref{sec:nocompose}).
  \item Only Survey and Admin are published, each through the reverse proxy
        (Section~\ref{sec:security-exposure}).
  \item The proxy refuses \opt{/api} on Survey, except \opt{/api/pdf/download}, and \opt{/q} on
        Survey and Admin (Section~\ref{sec:security-exposure}).
  \item Nothing but the proxy reaches an application's port, and the proxy replaces forwarding
        headers sent by a client (Section~\ref{sec:security-exposure}).
  \item \opt{ELICIT_TRUSTED_PROXIES} on Survey and Admin names the reverse proxy's address or
        network (Section~\ref{sec:security-exposure}).
  \item The proxy limits the rate of requests from one client address
        (Section~\ref{sec:security-exposure}).
  \item Report services, Pedigree, the database, the mail relay and the trace collector answer
        on the private network only (Table~\ref{tab:exposure}).
\end{itemize}

\subsubsection*{Transport}
\begin{itemize}[label = \elicitcheck]
  \item Survey and Admin are served over \val{https} only, and \val{http} redirects to it
        (Section~\ref{sec:security-transport}).
  \item A response from Survey and from Admin, fetched through the proxy, carries the headers
        of Table~\ref{tab:security-headers}.
  \item The issuer address and every invitation link are \val{https}
        (Section~\ref{sec:security-transport}).
  \item \opt{quarkus.tls.trust-all} is \val{false} on every module
        (Section~\ref{sec:security-transport}).
  \item Mail and the database connection are encrypted wherever they cross a network the site
        does not control (Section~\ref{sec:security-transport}).
\end{itemize}

\subsubsection*{Identity and accounts}
\begin{itemize}[label = \elicitcheck]
  \item The console signs in against the site's own provider, with an exact redirect address
        (Section~\ref{sec:security-identity}).
  \item The console's client allows the authorization code flow only
        (Section~\ref{sec:security-identity}).
  \item Every \val{elicit_admin} account has a second factor
        (Section~\ref{sec:security-identity}).
  \item The provider locks an account out after repeated failed sign-ins, and records
        sign-ins and administrative changes (Section~\ref{sec:security-identity}).
  \item \val{elicit_importer} is held by a machine client only and with no other Elicit
        role, and that client has an import client record with only the departments it
        registers into
        (Section~\ref{sec:security-identity}).
  \item The seeded accounts \val{admin} and \val{user} are renamed, and the console shows no
        banner about them (Chapter~\ref{sec:firstsignin}).
  \item Each user has only the departments they work in
        (Section~\ref{sec:security-identity}).
\end{itemize}

\subsubsection*{Access codes}
\begin{itemize}[label = \elicitcheck]
  \item \opt{accessCode.autoRegister} is \val{true} on Survey only where anyone may take the
        survey (Section~\ref{sec:security-access-codes}).
  \item The proxy's access log leaves out the path of invitation links, or the people who can
        read it are the people who may see access codes
        (Section~\ref{sec:security-observability}).
\end{itemize}

\subsubsection*{Database and secrets}
\begin{itemize}[label = \elicitcheck]
  \item No password in Table~\ref{tab:secrets} is in use anywhere on the site
        (Chapter~\ref{sec:secrets}).
  \item Every secret reaches its module from an environment variable or a secret store, and
        none is in a file under version control (Section~\ref{sec:security-secrets}).
  \item Every module holds the same service token, generated for this site; no module's log
        warns of the development value; and the trusted origins name only services the site
        runs (Chapter~\ref{sec:service-token}).
  \item Neither \val{elicit_owner} nor \val{dbowner} is a superuser
        (Section~\ref{sec:security-database}).
  \item A report service of the site's own connects as \val{surveyreport_user}
        (Section~\ref{sec:security-database}).
  \item Backups are encrypted and access to them is restricted
        (Section~\ref{sec:security-database}).
\end{itemize}

\subsubsection*{What the site loads and calls}
\begin{itemize}[label = \elicitcheck]
  \item Somebody is named as the person who decides which definition files are applied, and
        reads the addresses each one brings (Section~\ref{sec:security-content}).
  \item The brand and translations directories are mounted read-only and writable by operators
        only (Section~\ref{sec:security-content}).
  \item Every report address names a service on the private network, and every post-survey
        action that leaves the site is \val{https} (Section~\ref{sec:security-reports}).
  \item FHHS verifies the SFTP server's host key (Section~\ref{sec:security-reports}).
\end{itemize}

\subsubsection*{Operations}
\begin{itemize}[label = \elicitcheck]
  \item Logs and traces go only where personal data may go
        (Section~\ref{sec:security-observability}).
  \item Every module runs the same named release (Section~\ref{sec:security-releases}).
  \item Somebody is named as the person who watches for new releases and applies them
        (Section~\ref{sec:security-releases}).
\end{itemize}

%% ================================================================== 16
\section{Verifying the installation}
\label{sec:verify}

Work through these checks in order: each one relies on the one before it.

\begin{enumerate}
  \item \textbf{Every service is ready.} \file{./status.sh} lists the containers and asks each
        application whether it is ready: \opt{/q/health/ready} for Survey, Admin, FHHS,
        Author and Author's preview, \opt{/healthz} for Pedigree, and the front page of the
        bundled Keycloak and Mailpit. A service that is up but not ready is reported as
        such, with the reason its readiness check gives. Under your own container platform,
        ask each \opt{/q/health/ready} directly. For Survey, ready means its migrations have
        finished.
  \item \textbf{The modules authenticate each other.} From a host on the private network,
        ask each module directly, without the service token. On the evaluation stack the
        addresses are \val{localhost:8080}, \val{localhost:8082} and \val{localhost:8083}.
\begin{elicitcmd}[fontsize = \footnotesize]
curl -s -o /dev/null -w '%{http_code}\n' -X POST http://<survey>/api/etl/build
curl -s -o /dev/null -w '%{http_code}\n' -X POST \
     -H 'Authorization: Bearer not-this-token' http://<survey>/api/etl/build
curl -s -o /dev/null -w '%{http_code}\n' -X POST \
     -H 'Content-Type: application/json' -d '{}' http://<fhhs>/proband/report
curl -s -o /dev/null -w '%{http_code}\n' -X POST http://<pedigree>/svg
curl -s -o /dev/null -w '%{http_code}\n' http://<survey>/q/health/ready
curl -s -o /dev/null -w '%{http_code}\n' http://<pedigree>/healthz
\end{elicitcmd}
        The first four print \val{401} and the last two \val{200}. A \val{503} from one of
        the first four means that module has no token (Chapter~\ref{sec:service-token}). The
        other half needs no command: a definition that applies with its reporting schema
        rebuilt, and a report page that shows its reports, are calls that carried the token.
  \item \textbf{The schema exists.} The \val{survey} schema carries tables, and each module's
        migration history is present.
  \item \textbf{The reporting schema is built.} After a definition is applied, the survey's
        \val{report_schema} is set and \val{dim_step} in that schema has rows
        (Chapter~\ref{sec:definition}).
  \item \textbf{Sign-in works.} Sign in to Admin. A redirect that fails here is almost always
        the identity provider's issuer address (Chapter~\ref{sec:oidc}).
  \item \textbf{The console agrees.} Admin's \ui{System} section reports on the database,
        branding, mail, the reachability of every outbound dependency (report services
        included) and the current sign-in's roles. Its screens are shown below.
  \item \textbf{A respondent can answer.} Register a subject, issue an access code, and walk
        the survey to the end. If the survey has reports, the report page shows each of them.
  \item \textbf{The site is secured.} Work through the checklist that closes
        Chapter~\ref{sec:security} (Section~\ref{sec:security-checklist}). The checks above
        show that the site works; the checklist shows that it works only for the people it
        should.
\end{enumerate}

\screenshot{11-system-overview}{The System overview page. It lists the setup work still
outstanding, which on a finished installation makes it the shortest possible verification.}

\screenshot{12-system-database}{The System database page: both connections, and whether each
module's migrations are in place.}

\screenshot{14-system-email}{The System email page: the effective settings, and a button that
sends a real message through them.}

\screenshot{15-system-connections}{The System connections page: a bounded reachability check of
every outbound dependency, each of the survey's report services among them.}

\screenshot{16-system-oidc}{The System OIDC page: the current sign-in's roles, and its tokens
masked. Secrets are reported as present or absent, never shown.}

\begin{elicitnote}[Survey writes a diagnostics report to its log]
A few seconds after startup Survey writes a startup diagnostics report, which also checks
every report service it will call. Anything an operator must act on is logged as a warning;
the rest is logged at information level, and its category is pinned so the report appears even
when the overall log level is \val{WARN}. It is the first place to look after a start that
seemed to work.
\end{elicitnote}

%% ================================================================== 17
\section{When it does not work}
\label{sec:trouble}

Find the symptom in Table~\ref{tab:failures}: each row is a failure that has actually been seen
on an installation, with its cause and its fix. Where no row matches, read these, in order:
\begin{enumerate}
  \item the failing service's log, for the exception;
  \item Survey's startup diagnostics, for what it could and could not reach when it started;
  \item Admin's \ui{System} pages, for the state of the running site
        (Chapter~\ref{sec:verify});
  \item the database's migration history (\val{flyway_schema_history} in each schema), for
        which module reached which version.
\end{enumerate}

\begin{plaintable}{tab:failures}{Failures that are actually seen}{@{}>{\raggedright\arraybackslash}p{0.29\linewidth} >{\raggedright\arraybackslash}p{0.31\linewidth} >{\raggedright\arraybackslash}p{0.33\linewidth}@{}}{\thead{Symptom} & \thead{Cause} & \thead{Fix}}
A module fails its migration on a missing table or sequence. & It started before Survey
migrated a new database. & Let Survey become ready, then start the module again. Nothing is
corrupted. \\
\addlinespace
FHHS is not ready and answers report requests with \val{503}, naming a survey. & Its survey
has not been imported yet, or its reporting schema has not been built. & Apply the definition,
or apply it again to build the schema (Chapter~\ref{sec:definition}). FHHS goes healthy on its
next probe; nothing needs restarting. \\
\addlinespace
A module is not ready, and its \val{service-token} check gives the reason that
\opt{ELICIT_SERVICE_TOKEN} is not set. Its protected paths answer \val{503}. & The module was
started without the service token. & Give it the site's token and restart it
(Chapter~\ref{sec:service-token}). \\
\addlinespace
A module exits at startup, and its log says \opt{ELICIT_SERVICE_TOKEN} is too short. & The
token has fewer than 32 characters. & Generate one with \val{openssl rand -hex 32} and give
it to every module. \\
\addlinespace
A module exits at startup, and its log says an entry of
\opt{ELICIT_SERVICE_TRUSTED_ORIGINS} is not an origin. & The entry carries a path, a query or
a user name, or lacks a scheme or a host. & Write it as scheme, host and port, such as
\val{http://fhhs.internal:8080}. \\
\addlinespace
A module's log warns at every start that \opt{ELICIT_SERVICE_TOKEN} is the well-known
development value. & The module runs with the token the compose file ships. & Expected on an
evaluation stack. On a site, generate a token of its own
(Chapter~\ref{sec:service-token}). \\
\addlinespace
An apply reports that the reporting schema was not rebuilt, a report card shows an error
although its service is running, or FHHS's post-survey action fails, and the module that was
called answered \val{401}. & Either the two modules hold different service tokens, as they do
part-way through a rotation; or the address is at an origin the caller does not list as
trusted, so it sent no token. & Give every module the same token and restart them together;
or list the origin on Survey, Admin and FHHS, or change the addresses in Author to an origin
that is listed (Section~\ref{sec:service-token-origins}). The caller's log shows the call
failing with \val{401}. \\
\addlinespace
A report card shows an error, or a report service is unreachable on \ui{Connections}. & The
address in the survey's report list does not resolve from inside the Survey and Admin
containers, or the service is down. & Correct the address in Author and apply the new export
(Chapter~\ref{sec:reports}); a fix in the database lasts only until the next apply. \\
\addlinespace
A report shows on the respondent's page but is missing from the console's printout. & The
service answers no \val{pdf}. & Return the printable description as well
(Chapter~\ref{sec:reports}). \\
\addlinespace
A post-survey action's deliveries stay \val{CALL_REFUSED}. & The action sends the External ID
to an \val{http} address, uses a placeholder it may not, or needs an External ID the subject
lacks; the delivery's \val{error_msg} says which. & Correct the action in Author and apply the new
export, or give the subject an External ID. A refused delivery is not retried
(Section~\ref{sec:psa}). \\
\addlinespace
A post-survey action's calls fail with \val{401} or \val{403}. & No token reached the
service, or the service no longer accepts it: the setting's name does not hold the action's key,
or the token has expired or been revoked. & Check the line for the action in Survey's startup
diagnostics, which says whether it has a credential; if it does, ask the service for a new token
(Section~\ref{sec:psa}). \\
\addlinespace
Sign-in redirects fail, or land on the wrong service. & The issuer address does not resolve
to the same provider from the browser and from inside the container. & Use one address that
works from both (Chapter~\ref{sec:oidc}). \\
\addlinespace
A service is reported unhealthy although it responds in a browser. & The health check was
pointed at the published host port. & Point it at \val{8080} inside the container. \\
\addlinespace
The console shows a banner about default accounts on every screen. & The seeded accounts
\val{admin} and \val{user} have not been renamed. & Rename both
(Chapter~\ref{sec:firstsignin}). \\
\addlinespace
A dialog blocks the console and will not close. & No department exists. & Follow its link and
create one (Chapter~\ref{sec:firstsignin}). \\
\addlinespace
A signed-in user can only log out. & They belong to no department. & Assign them one in
\ui{Users}. \\
\addlinespace
A survey is missing from \ui{Register Subjects}, or a CSV row or API record is refused as not
permitted for its survey. & The subject's department does not list that survey; a newly applied
survey is in no department. & Add the survey to the department on \ui{Departments}
(Chapter~\ref{sec:definition}). \\
\addlinespace
Sign-in succeeds, but the console says the account has no active application record. & The
identity provider knows the account; the console has no user record for it. It is a missing
account, not a missing department. & Sign in with a seeded account and create the record on
\ui{Users}, with a department (Chapter~\ref{sec:firstsignin}). \\
\addlinespace
The applications still show the default brand. Admin's \ui{Branding} page reports every asset
as embedded. & The configured path does not match the volume's target inside the container, or
the container was restarted rather than recreated after the volume changed. & Compare
\opt{brand.file.system.path} with the mount, and recreate the container
(\val{docker compose up -d}), not restart it (Chapter~\ref{sec:brand}). \\
\addlinespace
The pedigree card shows an error in place of the drawing. & FHHS could not reach Pedigree in
the short time it waits, or Pedigree answered with an error, which the card shows. & Check that
Pedigree is running and that \opt{pedigree.url} reaches it from inside the FHHS container
(Chapter~\ref{sec:fhhs}). \\
\addlinespace
The database warns of a collation version mismatch, or creating a database fails on one. & The
data directory was made under another operating system or image, with a different C library. &
As a superuser, in each database the warning names, \val{REINDEX DATABASE <name>;} then
\val{ALTER DATABASE <name> REFRESH COLLATION VERSION;}, and repeat what failed. \\
\addlinespace
A mounted language is never offered. & Any one of these:
\begin{itemize}
  \item The directory is not mounted on that application, or \opt{i18n.file.system.path} does
        not name the mount point.
  \item The file is in the other application's sub-directory, or its name uses a hyphen
        (\file{translations_es-419.properties}).
  \item The tag is missing from \opt{i18n.offered.locales}.
  \item The module has not been restarted since the change.
\end{itemize}
& The application logs a directory it found nothing in, and each listed tag with no file. Fix
the cause (Section~\ref{sec:language}) and restart. \\
\end{plaintable}

\subsection{Upgrading an existing site}

Upgrade procedures belong to a release, and this manual, stamped with one version, may outlive
it. They are kept in \file{DeploymentScript.md} in the umbrella repository, which covers,
among others, the move to Kimball Type 2 slowly changing dimensions and the rename of the
respondent credential to the access code. Follow the procedure for the release you are moving
to, not this manual.

\begin{elicitcaution}[Upgrade the modules that share a schema together]
Survey, Admin and FHHS share the \val{survey} schema and migrate on top of each other.
Upgrading one without the others is what breaks. In the compose file, \opt{depends_on} orders
\emph{startup} only; it does not check again a dependency that is already running.
\end{elicitcaution}

%% ================================================================== 18
\section{Where to look next}

This manual stops once a site is installed, secured and verified. These carry on from it:

\begin{description}
  \item[The administrator's manual] Running the console: departments, users, subjects, access
        codes, invitations, reports, and applying survey definitions.
  \item[The Author installation manual] Provisioning the authoring tool: its database, its
        identity-provider client and roles, and the preview Survey it runs beside.
  \item[\file{DeploymentScript.md}] Release-specific upgrade procedures, and the post-install
        template data for the Family Health History Survey.
  \item[The FHHS repository] The source of the three Family History reports, one package per
        report, and its use cases under \file{docs/use_cases/}: the most complete example of
        a report service there is.
  \item[\file{docs/BRAND_SYSTEM_IMPLEMENTATION_GUIDE.md}] How the brand system is built, for
        someone extending it rather than mounting one.
  \item[\file{docs/I18N_IMPLEMENTATION_GUIDE.md}] The same, for translations.
  \item[\url{https://quarkus.io/guides/all-config}] Every configuration setting this manual
        does not list.
\end{description}

\vspace{1em}

\begin{elicitnote}[Which version is this?]
This manual describes the Elicit platform \textbf{\elicitVersion}, built
\textbf{\elicitBuildDate}. The same pair is printed at the foot of every page, so a printed
copy can always be matched back to the release it came from.
\end{elicitnote}

\elicitpart{Reference}{Every setting Survey and Admin define, every Quarkus setting their images
set, and how a setting reaches an application.}
\appendix
\renewcommand{\thesection}{\Alph{section}}

%% ================================================================== A
\section{Configuration reference}
\label{sec:config}

This appendix lists every setting Survey and Admin define and every Quarkus setting their
images set. FHHS's own settings are with its installation, in Table~\ref{tab:fhhs}; it shares
the datasource and observability settings below.

\subsection{How configuration reaches an application}

Every module is a Quarkus application, so every setting can be given as a property or as an
environment variable, whichever your deployment makes easier. When a setting is given in more
than one place, the source with the highest ordinal wins:

\begin{plaintable}{tab:sources}{Where a setting can come from}{@{}>{\raggedright\arraybackslash}p{0.62\linewidth} >{\raggedleft\arraybackslash}p{0.12\linewidth}@{}}{\thead{Source} & \thead{Ordinal}}
System properties, passed on the command line & 400 \\
Environment variables & 300 \\
Environment variables from a \file{.env} file in the working directory & 295 \\
\file{application.properties} in a \file{config} directory, relative to the application's
working directory in its container & 260 \\
The \file{application.properties} inside the image: the defaults this appendix lists & 250 \\
\end{plaintable}

A \file{.env} file therefore sits \emph{below} real environment variables, not above them.

\subsubsection*{From a property name to an environment variable}

Upper-case the name and replace every character that is neither a letter, a digit nor an
underscore with an underscore. So \opt{quarkus.datasource.jdbc.url} becomes
\opt{QUARKUS_DATASOURCE_JDBC_URL}, and a quoted segment produces a doubled underscore on each
side: \opt{quarkus.log.category."io.agroal".level} becomes
\opt{QUARKUS_LOG_CATEGORY__IO_AGROAL__LEVEL}.

\begin{elicitcaution}[The mapping is lossy]
Names that differ only in punctuation collide: \opt{foo.bar-baz} and \opt{foo.bar.baz} both
become \opt{FOO_BAR_BAZ}. Where that matters, give the setting as a property instead. In a
container you can also give the dotted name directly as an environment entry, which avoids
the question entirely and is what the shipped compose file does.
\end{elicitcaution}

\subsubsection*{Profiles}

A name prefixed with a profile, such as \optlit{\%dev.quarkus.http.port}, applies only under that
profile. Three prefixes appear in Survey's and Admin's files, and none of them configures a
deployment:

\begin{itemize}
  \item \optlit{\%dev.}: values for \val{quarkus:dev} on a developer's machine.
  \item \optlit{\%test.}: values for the module's own test run.
  \item \optlit{\%docker.}: a build-time profile that does one thing, switch the container-image
        build on. It sets nothing a running deployment reads.
\end{itemize}

There are no \optlit{\%prod.} entries. Everything a deployment runs with is either an unprefixed
default from the tables below or a value your site supplies.

\subsection{Reading the tables}

The \textbf{Change} column says what the setting is to you:

\begin{description}
  \item[\chgset] Your site must set it. The shipped value is a development value or a
        placeholder and is wrong anywhere else.
  \item[\chgmay] Your site may set it. The shipped value is a working default.
  \item[\chgkeep] Leave it alone. It is an internal default that happens to live in the same
        file, and changing it changes how the application works.
\end{description}

\subsection{Settings common to Survey and Admin}

\begin{settings}{tab:common}{Settings Survey and Admin both read}
\opt{brand.file.system.path} & \val{/brand} & \chgmay & The mounted brand directory. The
compose file's commented brand mount uses \val{/opt/brand}. See Chapter~\ref{sec:brand}. \\
\opt{brand.local.path} & \val{brand} & \chgkeep & A brand directory beside the running
application, searched between the mount and the embedded default. \\
\opt{i18n.file.system.path} & \val{/i18n} & \chgmay & The mounted translations directory. Each
application reads its own sub-directory of it (\opt{i18n.app.name}). The compose file's commented
translations mount puts \file{elicit-i18n/} at \val{/opt/i18n}. A directory that does not exist is
logged and the application offers English only. See Chapter~\ref{sec:i18n}. \\
\opt{i18n.app.name} & per app & \chgkeep & This application's sub-directory of the translations
directory: \val{survey} or \val{admin}. Must be a plain directory name. \\
\opt{i18n.offered.locales} & --- & \chgmay & The languages this application offers, when it
should offer fewer than are mounted: a comma-separated list of tags. Unset offers every mounted
language. English is always offered. A listed tag with no mounted file is logged and not
offered, so this cannot add a language. A language left out is hidden from the selector, a
\opt{?lang=} naming it is refused, and survey content is not served in it either, because
content is only offered in a language the site also offers for the chrome. See
Chapter~\ref{sec:i18n}. \\
\optlit{i18n.direction.<tag>} & --- & \chgmay & Overrides both manifests (the image's
\file{META-INF/i18n/i18n-config.json} and the \file{i18n-config.json} at the root of the
translations directory) for one language tag: \val{rtl} or \val{ltr}. Both manifests declare
\val{ar} right-to-left; every other language is inferred from its tag. Any other value is logged
and ignored. See Section~\ref{sec:direction}. \\
\optlit{i18n.font-scale.<tag>} & --- & \chgmay & Overrides both manifests' \val{fontScale} for
one tag, multiplying the root font size on pages shown in that language. Both manifests scale
\val{ar} to \val{1.15}; every other language is \val{1.0}. Must be a number from \val{0.75} to
\val{2.0}; a value outside that range is logged and dropped rather than clamped. See
Section~\ref{sec:direction}. \\
\opt{i18n.pseudo-locale.enabled} & \val{false} & \chgkeep & Renders every text as a marker.
A test aid; never enable it on a site. \\
\opt{quarkus.application.name} & per app & \chgkeep & The name this instance reports in
metrics and traces. \\
\opt{build.timestamp} & build & \chgkeep & Stamped by the build and shown on the
\ui{System} overview. \\
\end{settings}

\subsection{Survey}

\begin{settings}{tab:survey}{Survey: Elicit settings}
\opt{accessCode.autoRegister} & \val{false} & \chgmay & Accept \emph{any} access code of 5 to 12
characters rather than only one the console issued, creating a respondent with no subject
(Section~\ref{sec:security-access-codes}). Set \val{true} for a survey anyone may take; the
compose file sets it for evaluation. \\
\opt{elicit.etl.enabled} & \val{true} & \chgmay & Whether this instance builds the reporting
star schema of every installed survey at startup and updates it when a respondent finalizes.
Set \val{false} only on an instance that merely renders surveys and that nothing reports
on. \\
\opt{elicit.etl.drop.enabled} & \val{false} & \chgmay & Whether \val{DELETE
/api/etl/schema/<key>} may drop a survey's reporting schema. Leave it \val{false} at a site:
the next build would regenerate the schema with new surrogate ids, which breaks any extract
keyed on them. Only a database that holds no reporting anyone depends on turns it on. \\
\opt{elicit.service.token} & --- & \chgset & The site's service token: the secret every module
presents when it calls another, and that Survey requires on \opt{/api}
(Chapter~\ref{sec:service-token}). At least 32 characters; no value is shipped. A secret.
Environment: \opt{ELICIT_SERVICE_TOKEN}. \\
\opt{elicit.service.trusted-origins} & --- & \chgmay & The origins of the report services and
post-survey actions Survey may send the service token to: a comma-separated list such as
\val{http://fhhs.internal:8080}. Unset, Survey sends the token to none
(Section~\ref{sec:service-token-origins}). Environment:
\opt{ELICIT_SERVICE_TRUSTED_ORIGINS}. \\
\opt{elicit.diagnostics.startup.enabled} & \val{true} & \chgmay & Whether the startup
diagnostics report is written to the log. \\
\opt{elicit.diagnostics.startup.delay} & \val{10s} & \chgmay & How long after startup that
report is written. \\
\opt{elicit.content.translations-ttl} & \val{PT5M} & \chgmay & How long a survey's content
translations are held in memory before being reloaded, as an ISO-8601 duration. A translation
published through \emph{Apply Survey Definition} reaches respondents within this window without
a restart. See Chapter~\ref{sec:i18n}. \\
\opt{elicit.content.serve-stale-translations} & \val{false} & \chgmay & Whether to show a
translation whose base text has been edited since it was made. Left \val{false}, such a string
falls back to the survey's base language, so every language asks the same question; set
\val{true} only where continuity matters more than that. \\
\opt{elicit.post-survey-action.retry.enabled} & \val{true} & \chgmay & Whether a post-survey
action whose call failed is called again on a schedule. See Section~\ref{sec:psa}. \\
\opt{elicit.post-survey-action.retry.interval} & \val{15m} & \chgmay & How often the retry
runs, as a duration (\val{15m}, \val{1h}). \\
\opt{elicit.post-survey-action.retry.max-attempts} & \val{24} & \chgmay & How many attempts a
delivery gets, the first call included, before Survey stops retrying it. \\
\end{settings}

A post-survey action's credential is set per action, under a name that holds the action's key:
\optlit{elicit.post-survey-action."<key>".bearer-token} (Section~\ref{sec:psa}). It is a
secret and has no default.

\subsection{Admin}

\begin{settings}{tab:admin}{Admin: Elicit settings}
\opt{elicit.authorization.mode} & \val{OIDC} & \chgmay & \val{DATABASE} also consults
\val{survey.user_roles} and shows the role-assignment section of the user editor. It does not
disable OIDC. Environment: \opt{ELICIT_AUTHORIZATION}. \\
\opt{elicit.survey.url} & \val{http://survey:8080} & \chgset & Where Survey is, for the
reporting-schema rebuild after a definition is applied and for the respondent reports Admin
prints. The default is the compose service
name. Admin sends the service token to this address. Environment: \opt{ELICIT_SURVEY_URL}. \\
\opt{elicit.service.token} & --- & \chgset & The site's service token, which Admin presents to
Survey and to report services at a trusted origin (Chapter~\ref{sec:service-token}). At least
32 characters; no value is shipped. A secret. Environment: \opt{ELICIT_SERVICE_TOKEN}. \\
\opt{elicit.service.trusted-origins} & --- & \chgmay & The origins of the report services Admin
may send the service token to, besides Survey's: a comma-separated list such as
\val{http://fhhs.internal:8080}. Unset, Admin sends the token to Survey only
(Section~\ref{sec:service-token-origins}). Environment:
\opt{ELICIT_SERVICE_TRUSTED_ORIGINS}. \\
\opt{elicit.survey.etl-build.enabled} & \val{true} & \chgmay & Whether to make that call at
all. Environment: \opt{ELICIT_SURVEY_ETL_BUILD_ENABLED}. \\
\opt{elicit.mailer.send-timeout-seconds} & \val{30} & \chgmay & How long a single message
send may take before it is abandoned. \\
\opt{elicit.participation.small-cell} & \val{5} & \chgmay & On the Participation page and in its
CSV and PDF, a count from 1 to one below this value is shown as ``< \textit{k}'' (\textit{k} being this value) and the rates built on it
are withheld, so that a narrow filter cannot single out a respondent. \val{1} turns suppression
off. Environment: \opt{ELICIT_PARTICIPATION_SMALL_CELL}. \\
\opt{elicit.participation.stalled-days} & \val{14} & \chgmay & On the Participation page and in its
exports, a respondent in progress who has not opened a page or saved an answer for this many days
is counted as stalled. Environment: \opt{ELICIT_PARTICIPATION_STALLED_DAYS}. \\
\opt{elicit.debug.reveal-tokens} & \val{false} & \chgkeep & Shows OIDC and bearer tokens
unmasked on the diagnostics screen, in development mode only: a released image ignores it, and
\ui{System} says so when it is set. This concerns identity tokens, not respondent access
codes. \\
\opt{admin.manual.resource} & \val{manual/elicit-admin-manual.pdf} & \chgkeep & Where the
administrator's manual sits on the classpath. A build carrying no manual simply does not
offer it. \\
\opt{quarkus.mailer.mock} & --- & \chgkeep & When true, messages are logged instead of sent.
Set by Quarkus itself in the development and test profiles; never set it on a site. \\
\end{settings}

\subsection{The Quarkus settings the images set}

These are not Elicit settings, but the images set them, so your site is overriding a value
rather than a blank when it changes one.

\begin{settings}{tab:quarkus-db}{Datasources: Survey, Admin and FHHS}
\opt{quarkus.datasource.db-kind} & \val{postgresql} & \chgkeep & Elicit runs on PostgreSQL
only. \\
\opt{quarkus.datasource.db-version} & \val{17} & \chgkeep & The version the schema and the
migrations are written against. \\
\opt{quarkus.datasource.jdbc.url} & \val{jdbc:postgresql://db:5432/survey} plus the
parameters below & \chgset & The application connection. The host names the compose
service. \\
\opt{quarkus.datasource.username} & per app & \chgset & \val{survey_user} for Survey and
FHHS, \val{surveyadmin_user} for Admin. \\
\opt{quarkus.datasource.password} & none & \chgset & No value is shipped. Environment:
\opt{ELICIT_DB_PASSWORD}. See Chapter~\ref{sec:secrets}. \\
\opt{quarkus.datasource.jdbc.min-size} & \val{5} & \chgmay & \\
\opt{quarkus.datasource.jdbc.max-size} & \val{20} & \chgmay & \\
\opt{quarkus.datasource.jdbc.leak-detection-interval} & \val{10m} & \chgmay & \\
\opt{quarkus.datasource.owner.jdbc.url} & as above & \chgset & The migrating connection, as
which every module's Flyway migrations run. \\
\opt{quarkus.datasource.owner.username} & \val{elicit_owner} & \chgkeep & The schema owner. \\
\opt{quarkus.datasource.owner.password} & none & \chgset & No value is shipped.
Environment: \opt{ELICIT_OWNER_DB_PASSWORD}. See Chapter~\ref{sec:secrets}. \\
\end{settings}

The shipped connection URLs carry four parameters that the host and database name above
leave out. Keep them when you point a module at your own server: they are what stops a
dropped connection hanging a request indefinitely:

\begin{elicitfile}[fontsize=\footnotesize]
?connectTimeout=30&socketTimeout=600&tcpKeepAlive=true&tcpKeepAlivesIdle=60
\end{elicitfile}

When the database is on another host, encrypt the connection and verify the server by adding
two more to both URLs, the second naming the certificate authority that signed the server's
certificate (Section~\ref{sec:security-transport}):

\begin{elicitfile}[fontsize=\footnotesize]
&sslmode=verify-full&sslrootcert=<path to the CA certificate>
\end{elicitfile}

\begin{settings}{tab:quarkus-oidc}{Identity: Admin only}
\opt{quarkus.oidc.auth-server-url} & \val{http://host.docker.internal:8180/realms/elicit/} &
\chgset & Your provider's issuer address: \val{https} on a real site
(Section~\ref{sec:security-transport}). It must resolve to the same provider from the
browser and from inside the container. \\
\opt{quarkus.oidc.client-id} & \val{elicit-admin} & \chgmay & \\
\opt{quarkus.oidc.credentials.secret} & \opt{OIDC_CLIENT_SECRET} & \chgset & The
client secret. \\
\opt{quarkus.oidc.application-type} & \val{hybrid} & \chgkeep & Accepts browser sign-ins and
bearer tokens on one client. \\
\opt{quarkus.oidc.authentication.redirect-path} & \val{/callback} & \chgkeep & \\
\opt{quarkus.oidc.logout.path} & \val{/logout} & \chgkeep & \\
\end{settings}

\begin{settings}{tab:quarkus-http}{HTTP, proxying and limits}
\opt{quarkus.http.port} & \val{8080} & \chgkeep & Inside the container. Publish a different
host port instead. \\
\opt{quarkus.http.proxy.proxy-address-forwarding} & \val{true} & \chgkeep & Required behind a
reverse proxy or load balancer so that redirect URIs are built with the external address. Safe
when only the proxy reaches the application's port
(Section~\ref{sec:security-exposure}). \\
\opt{quarkus.http.proxy.allow-x-forwarded} & \val{true} & \chgkeep & Honors the
\val{X-Forwarded-*} headers. Have the proxy set them itself and discard any a client
sent. \\
\opt{quarkus.http.proxy.trusted-proxies} & loopback and the private ranges & \chgset & The
addresses whose forwarding headers Survey and Admin honor; from any other address they are
ignored. The default is \val{127.0.0.1}, \val{::1}, \val{10.0.0.0/8}, \val{172.16.0.0/12},
\val{192.168.0.0/16} and \val{fc00::/7}. Narrow it to your reverse proxy's address or
network, comma-separated. Environment: \opt{ELICIT_TRUSTED_PROXIES}. \\
\opt{quarkus.http.cors.origins} & empty & \chgmay & Admin lists allowed origins explicitly;
Survey defaults to \val{http://localhost:8080}. Each application is served same-origin to
its own interface, so a site normally needs nothing here. \\
\opt{quarkus.http.limits.max-body-size} & \val{25M} & \chgmay & Admin only. Bounds a CSV or
definition-file upload. \\
\opt{quarkus.tls.trust-all} & \val{false} & \chgkeep & Survey and Admin verify the
certificate of every service they call over TLS. See the note below. \\
\end{settings}

\begin{elicitnote}[Services signed by a private certificate authority]
Survey and Admin verify every outbound TLS connection (to the identity provider, the mail
relay, report services and each other) against the Java runtime's trust store. Releases before 3.0.0
shipped \opt{quarkus.tls.trust-all}\val{=true}, which turned that verification off; a site
that carried the setting forward should remove it. A service whose certificate is issued by a
private authority is trusted by adding that authority, for example
\opt{quarkus.tls.trust-store.pem.certs} naming its PEM file, not by turning verification off.
\end{elicitnote}

\begin{settings}{tab:quarkus-flyway}{Migrations: Survey and Admin}
\opt{quarkus.flyway.owner.schemas} & \val{survey} & \chgkeep & \\
\opt{quarkus.flyway.owner.migrate-at-start} & per app & \chgkeep & \val{true} on Admin.
\val{false} on Survey, which runs its migrations itself so that it can choose between the
current schema and the upgrade track. \\
\opt{quarkus.flyway.owner.baseline-on-migrate} & \val{true} & \chgkeep & \\
\opt{quarkus.flyway.owner.validate-on-migrate} & \val{true} & \chgkeep & \\
\opt{quarkus.flyway.owner.repair-at-start} & \val{true} & \chgkeep & Admin only. Rewrites the
recorded checksums of two migrations that were edited after release, so existing databases
start without a manual repair. \\
\opt{quarkus.flyway.owner.table} & per app & \chgkeep & \val{flyway_history} for Survey,
\val{flyway_admin_history} for Admin: each module keeps its own history in the shared
schema. \\
\opt{quarkus.flyway.owner.connect-retries} & \val{10} & \chgmay & How many times to retry a
database that is not up yet. \\
\opt{quarkus.flyway.owner.locations} & \val{db/migration} & \chgkeep & Which migration scripts
to run. Survey reads this itself to choose between the current schema and the upgrade track;
changing it points a release at scripts it was not built with. \\
\opt{quarkus.flyway.owner.baseline-version} & \val{0} & \chgkeep & The version recorded when
an existing database is brought under Flyway's control. \\
\opt{quarkus.flyway.owner.baseline-description} & \val{Empty Database} & \chgkeep & The
description recorded alongside it. \\
\end{settings}

\begin{settings}{tab:quarkus-mail}{Mail: Admin only}
\opt{quarkus.mailer.host} & \val{localhost} & \chgset & Your SMTP relay. \\
\opt{quarkus.mailer.port} & \val{25} & \chgset & \\
\opt{quarkus.mailer.from} & --- & \chgset & The address invitations are sent from. \\
\opt{quarkus.mailer.tls} & \val{false} & \chgset & \val{true} for a relay reached across a
network the site does not control; invitations carry access codes
(Section~\ref{sec:security-transport}). \val{false} suits the bundled mail catcher. \\
\opt{quarkus.mailer.username} & --- & \chgmay & With
\opt{quarkus.mailer.password} if your relay authenticates. \\
\end{settings}

\begin{settings}{tab:quarkus-obs}{Health, metrics and traces}
\opt{quarkus.smallrye-health.root-path} & \val{/q/health} & \chgkeep & Gives
\opt{/q/health/ready} and \opt{/q/health/live}. On Admin they answer without a sign-in, for
your container platform's probes: have the reverse proxy refuse \opt{/q} and read them on the
private network (Section~\ref{sec:security-observability}). \\
\opt{quarkus.micrometer.export.prometheus.path} & \val{/q/metrics} & \chgkeep & Answers
without a sign-in on Admin too, for a scraper your site provides on the private network. \\
\opt{quarkus.otel.enabled} & \val{true} & \chgmay & \\
\opt{quarkus.otel.exporter.otlp.endpoint} & \val{http://localhost:4317} & \chgset & Your
collector. The compose file sets \val{http://jaeger:4317}. \\
\opt{quarkus.otel.traces.sampler} & \val{parentbased_always_on} & \chgmay & Samples every
trace. On a busy site, \val{parentbased_traceidratio} with a ratio keeps a share of them. Traces
carry request addresses, so send them only to a collector the site protects
(Section~\ref{sec:security-observability}). \\
\opt{quarkus.log.level} & \val{WARN} & \chgmay & Environment: \opt{LOG_LEVEL}. FHHS
defaults to \val{INFO}. \\
\opt{quarkus.log.category."com.elicitsoftware".level} & \val{INFO} & \chgmay & Admin only.
Environment: \opt{ELICIT_LOG_LEVEL}. \\
\opt{quarkus.http.access-log.enabled} & \val{true} & \chgmay & Survey's log leaves out
invitation links and query strings, so it holds no access codes
(Section~\ref{sec:security-observability}). \\
\opt{quarkus.hibernate-orm.statistics} & \val{true} & \chgmay & Useful, and not free. \\
\end{settings}

\subsection{Everything else}

These tables cover what Elicit defines and what its images set. Every other Quarkus setting
behaves as Quarkus documents it:

\begin{center}
\url{https://quarkus.io/guides/all-config}
\end{center}

\end{document}
